\documentclass[12pt,a4paper]{report}

\usepackage[french]{babel}
\usepackage[utf8]{inputenc}
\usepackage[T1]{fontenc}
\usepackage{fancyhdr}
\usepackage{xcolor}
\usepackage{titlesec}
\usepackage{tikz}
\usepackage{setspace}
\usepackage[margin=2.5cm]{geometry}
\usepackage{tcolorbox}
\usepackage{amsmath}
\usepackage{graphicx}
\usepackage{longtable}
\usepackage{booktabs}
\usepackage{array}
\newcolumntype{R}[1]{>{\raggedleft\arraybackslash}p{#1}}
\newcolumntype{L}[1]{>{\raggedright\arraybackslash}p{#1}}
\usepackage{listings}
\usepackage{url}

\lstdefinestyle{pythoncode}{
  basicstyle=\ttfamily\footnotesize,
  backgroundcolor=\color{gray!10},
  frame=single,
  breaklines=true,
  columns=fullflexible,
  showstringspaces=false,
  keywordstyle=\bfseries,
  language=Python,
}
\lstset{style=pythoncode}

% Couleur du bandeau (violet, à ajuster selon la charte de l'établissement)
\definecolor{accentbar}{RGB}{124, 77, 160}

% L'en-tête doit toujours afficher le nom du CHAPITRE (pas la section
% courante) : \leftmark contient le nom de chapitre et n'est jamais écrasé
% par les \section, contrairement à \rightmark (mis à jour par chaque
% \section). On redéfinit \chaptermark pour qu'il stocke le titre du
% chapitre tel quel, sans le préfixe "Chapitre N." ni majuscules.
\renewcommand{\chaptermark}[1]{\markboth{#1}{#1}}

% En-tête de page : titre du chapitre à droite + petit carré violet, ligne fine en dessous
\pagestyle{fancy}
\fancyhf{}
\renewcommand{\headrulewidth}{0.4pt}
\renewcommand{\headrule}{\color{accentbar}\hrule height\headrulewidth}
\fancyhead[R]{\normalfont\itshape\leftmark\hspace{6pt}%
  \raisebox{-2pt}{\tikz\fill[accentbar] (0,0) rectangle (0.4cm,0.4cm);}}
\fancyfoot[R]{\colorbox{accentbar}{\textcolor{white}{\thepage}}}

% La classe "report" applique automatiquement le style "plain" (sans en-tête)
% à la première page de chaque \chapter — on le redéfinit pour qu'il soit
% identique au style "fancy" ci-dessus, afin que le bandeau violet apparaisse
% aussi sur ces pages de titre de chapitre.
\fancypagestyle{plain}{%
  \fancyhf{}
  \renewcommand{\headrulewidth}{0.4pt}
  \renewcommand{\headrule}{\color{accentbar}\hrule height\headrulewidth}
  \fancyhead[R]{\normalfont\itshape\leftmark\hspace{6pt}%
    \raisebox{-2pt}{\tikz\fill[accentbar] (0,0) rectangle (0.4cm,0.4cm);}}
  \fancyfoot[R]{\colorbox{accentbar}{\textcolor{white}{\thepage}}}
}

% Titre de type "Introduction générale" / "Résumé" / "Abstract" : centré, gras, sans numéro
\titleformat{\chapter}[display]
  {\normalfont\bfseries\Large\centering}{}{0pt}{}
\titlespacing*{\chapter}{0pt}{10pt}{25pt}

\setstretch{1.15}

\begin{document}

% ============================================================
% PAGE DE GARDE
% ============================================================
\begin{titlepage}
\thispagestyle{empty}

% Bandeau décoratif violet en diagonale (coin inférieur gauche)
\begin{tikzpicture}[remember picture, overlay]
  \fill[accentbar!25] ([shift={(0cm,-10cm)}]current page.north west) --
    ([shift={(6cm,-19cm)}]current page.north west) --
    ([shift={(0cm,-19cm)}]current page.north west) -- cycle;
  \fill[accentbar] ([shift={(0cm,0cm)}]current page.south west) --
    ([shift={(0cm,7cm)}]current page.south west) --
    ([shift={(9cm,0cm)}]current page.south west) -- cycle;
\end{tikzpicture}

\begin{minipage}[t]{0.15\textwidth}
  \centering
  \fbox{\parbox[c][1.5cm][c]{1.3cm}{\centering\tiny [LOGO\\ministère]}}
\end{minipage}
\hfill
\begin{minipage}[t]{0.55\textwidth}
  \centering\small
  Ministère de l'Enseignement Supérieur\\ et de la Recherche Scientifique
\end{minipage}
\hfill
\begin{minipage}[t]{0.2\textwidth}
  \raggedleft\small
  Année Universitaire\\ \textbf{2026/2027}
\end{minipage}

\vspace{1cm}

\begin{center}
  \fbox{\parbox[c][2cm][c]{3cm}{\centering\small [LOGO\\EPI Digital\\School]}}

  \vspace{0.3cm}
  \small EPI Digital School\\École Internationale Supérieure Privée Polytechnique de Sousse

  \vspace{2.5cm}
  {\LARGE\bfseries RAPPORT DE\\[4pt] STAGE PROFESSIONNEL}

  \vspace{0.5cm}
  Pour l'obtention du Diplôme National d'Ingénieur / de Mastère en

  \vspace{0.5cm}
  \textbf{Spécialité :} Parcours IA \& Data Science

  \vspace{1cm}
  {\Large\bfseries Intitulé}

  \vspace{0.5cm}
  \fcolorbox{accentbar}{white}{%
    \parbox{0.8\textwidth}{\centering\bfseries
    Segmentation automatique et suivi cinétique de plaies\\
    (scratch assay) par deep learning,\\
    avec interprétation assistée par LLM (RAG)}%
  }

  \vspace{2cm}
  \begin{minipage}{0.3\textwidth}\centering
    Lieu du stage\\ \textbf{Laboratoire de Biologie,\\Faculté de Médecine et de Pharmacie de Sousse}
  \end{minipage}
  \begin{minipage}{0.3\textwidth}\centering
    Réalisé par\\ \textbf{Karima Chiha}
  \end{minipage}
  \begin{minipage}{0.3\textwidth}\centering
    Encadré par\\ \textbf{[À COMPLÉTER]}
  \end{minipage}
\end{center}

\end{titlepage}

% ============================================================
% TABLE DES MATIÈRES
% ============================================================
\renewcommand{\contentsname}{Table des matières}
\tableofcontents
\thispagestyle{empty}
\newpage

% ============================================================
% LISTE DES FIGURES
% ============================================================
\listoffigures
\thispagestyle{empty}
\newpage

% ============================================================
% LISTE DES TABLEAUX
% ============================================================
\listoftables
\thispagestyle{empty}
\newpage

% ============================================================
% LISTE DES ABRÉVIATIONS
% ============================================================
\chapter*{Liste des abréviations}
\addcontentsline{toc}{chapter}{Liste des abréviations}
\markboth{Liste des abréviations}{Liste des abréviations}

\begin{longtable}{l p{10.5cm}}
\toprule
\textbf{Abréviation} & \textbf{Signification} \\
\midrule
\endfirsthead
\toprule
\textbf{Abréviation} & \textbf{Signification} \\
\midrule
\endhead
API & Application Programming Interface (interface de programmation applicative) \\
BCE & Binary Cross-Entropy (entropie croisée binaire) \\
BM25 & Best Matching 25 (algorithme de pondération lexicale pour la recherche d'information) \\
CD & Déploiement Continu (Continuous Deployment) \\
CI & Intégration Continue (Continuous Integration) \\
CLAHE & Contrast Limited Adaptive Histogram Equalization \\
CNN & Convolutional Neural Network (réseau de neurones convolutif) \\
CPU & Central Processing Unit \\
CRISP-DM & Cross-Industry Standard Process for Data Mining \\
GPU & Graphics Processing Unit \\
IA & Intelligence Artificielle \\
IC & Intervalle de Confiance \\
IoU & Intersection over Union \\
JPEG & Joint Photographic Experts Group \\
LLM & Large Language Model (grand modèle de langage) \\
MAE & Mean Absolute Error (erreur absolue moyenne) \\
MLOps & Machine Learning Operations \\
MRR & Mean Reciprocal Rank \\
RAG & Retrieval-Augmented Generation (génération augmentée par récupération) \\
RGB & Red Green Blue (rouge-vert-bleu) \\
TTA & Test-Time Augmentation (augmentation au moment du test) \\
\bottomrule
\end{longtable}

\newpage

% ============================================================
% RÉSUMÉ
% ============================================================
\chapter*{Résumé}
\addcontentsline{toc}{chapter}{Résumé}
\markboth{Résumé}{Résumé}

\noindent
\hspace{1cm}Le scratch assay est une méthode in vitro largement utilisée pour étudier la migration cellulaire et la fermeture de plaie, mais l'analyse quantitative des surfaces de plaie reste généralement manuelle, ce qui la rend longue et sujette à variabilité. Ce projet développe une solution basée sur l'intelligence artificielle pour la segmentation automatique et l'analyse quantitative d'images de microscopie de scratch assay.

\vspace{0.3cm}
\noindent
\hspace{1cm}Un jeu de 97 images de microscopie, annotées par des contours de plaie tracés à la main, a été transformé en masques binaires fiables via un pipeline de prétraitement dédié. Un modèle U-Net à encodeur ResNet34 pré-entraîné a été entraîné sur un split train/validation/test reproductible, puis évalué sur un jeu de test tenu à l'écart face à une baseline classique (analyse de texture et seuillage d'Otsu). Un ensemble de plusieurs modèles a également été étudié pour améliorer la robustesse et fournir un signal de confiance au niveau du pixel, fondé sur l'accord entre modèles. Une application Streamlit a été développée pour la segmentation automatique, la mesure, la visualisation et la correction manuelle, complétée par un assistant à génération augmentée par récupération (RAG) pour l'interprétation des résultats à la lumière de la littérature.

\vspace{0.3cm}
\noindent
\hspace{1cm}Sur le jeu de test, le U-Net a obtenu un score de Dice de 0,882 contre 0,576 pour la baseline, différence statistiquement significative (test de Wilcoxon, p = 0,000122). Un ensemble de quatre modèles a par ailleurs obtenu un Dice de 0,893 sur ce même jeu de test --- un résultat exploratoire, deux des quatre modèles ayant partagé des données d'entraînement avec cet échantillon ---, le désaccord entre modèles étant néanmoins corrélé négativement à la qualité de segmentation (r = -0,751) et offrant un signal de confiance supplémentaire. Ces résultats montrent que l'approche proposée surpasse la méthode classique testée, tout en reconnaissant explicitement ses limites et le besoin de validation complémentaire.

\vspace{0.3cm}
\noindent
\textbf{Mots-clés :} Scratch Assay, Deep Learning, Segmentation d'image, U-Net, Vision par ordinateur, Cicatrisation, Génération augmentée par récupération (RAG)

% ============================================================
% ABSTRACT
% ============================================================
\chapter*{Abstract}
\addcontentsline{toc}{chapter}{Abstract}
\markboth{Abstract}{Abstract}

\noindent
\hspace{1cm}The scratch assay is a widely used in-vitro method for studying cell migration and wound closure, but the quantitative analysis of wound areas is typically performed manually, making it time-consuming and subject to variability. This project develops an artificial intelligence-based solution for automated segmentation and quantitative analysis of scratch assay microscopy images.

\vspace{0.3cm}
\noindent
\hspace{1cm}A dataset of 97 microscopy images with manually drawn wound contours was processed into reliable binary masks through a dedicated pipeline. A U-Net model with a pretrained ResNet34 encoder was trained on a reproducible train/validation/test split and evaluated on a held-out test set against a classical image-processing baseline (texture analysis with Otsu thresholding). An ensemble of several models was also investigated to improve robustness and provide a pixel-level confidence signal based on model agreement. A Streamlit application was developed for automatic segmentation, measurement, visualization, and manual correction, complemented by a retrieval-augmented generation (RAG) assistant for literature-based interpretation of the results.

\vspace{0.3cm}
\noindent
\hspace{1cm}On the held-out test set, the U-Net achieved a Dice score of 0.882 versus 0.576 for the baseline, a statistically significant difference (Wilcoxon signed-rank test, p = 0.000122). A four-model ensemble further obtained a Dice score of 0.893 on this same test set --- an exploratory result, as two of the four models shared training data with this sample ---, with model disagreement correlating negatively with segmentation quality (r = -0.751) and offering an additional confidence signal. These results show that the proposed deep learning approach outperformed the tested classical method for wound segmentation, while its limitations and the need for further validation are explicitly acknowledged.

\vspace{0.3cm}
\noindent
\textbf{Keywords:} Scratch Assay, Deep Learning, Image Segmentation, U-Net, Computer Vision, Wound Healing, Retrieval-Augmented Generation (RAG)

% ============================================================
% INTRODUCTION GÉNÉRALE
% ============================================================
\chapter*{Introduction générale}
\addcontentsline{toc}{chapter}{Introduction générale}
\markboth{Introduction générale}{Introduction générale}

\noindent
\hspace{1cm}L'analyse de la migration cellulaire et de la fermeture des plaies constitue un aspect important de nombreuses études biologiques réalisées in vitro. Parmi les méthodes utilisées pour étudier ces phénomènes, le scratch assay, également appelé test de cicatrisation ou test de fermeture de plaie, est une technique couramment employée pour évaluer la capacité des cellules à migrer et à recouvrir progressivement une zone initialement dépourvue de cellules. Cette évolution peut être suivie à différents instants à partir d'images obtenues par microscopie, puis quantifiée à partir de la surface de la zone de plaie.

\vspace{0.3cm}
\noindent
\hspace{1cm}Cependant, l'analyse quantitative des images de scratch assay est souvent réalisée manuellement à l'aide d'outils de traitement d'images tels qu'ImageJ. Cette approche nécessite notamment de délimiter la zone de la plaie et de mesurer sa surface pour chaque image. Lorsque le nombre d'images augmente, cette opération devient longue et dépendante de l'intervention de l'utilisateur. Elle peut également introduire une certaine variabilité dans la délimitation des contours. L'automatisation de cette étape représente donc un intérêt pour faciliter l'analyse et améliorer sa reproductibilité.

\vspace{0.3cm}
\noindent
\hspace{1cm}Le choix de ce sujet s'inscrit dans la volonté d'appliquer les connaissances acquises en intelligence artificielle et en Data Science à une problématique concrète issue du domaine biologique. Ce projet offre ainsi l'opportunité de mettre en pratique différentes compétences en traitement d'images, Deep Learning, analyse de données et développement d'applications, tout en répondant à un besoin d'automatisation identifié dans le cadre du stage.

\vspace{0.3cm}
\noindent
\hspace{1cm}Dans ce contexte, l'intelligence artificielle et plus particulièrement les techniques de vision par ordinateur et de Deep Learning offrent des possibilités pour automatiser la segmentation des régions d'intérêt dans les images biologiques. Les réseaux de neurones dédiés à la segmentation, tels que U-Net, permettent notamment d'identifier automatiquement les pixels appartenant à une région donnée. Néanmoins, l'utilisation d'un modèle de Deep Learning ne garantit pas à elle seule des résultats fiables. Il est nécessaire d'évaluer ses performances sur des données indépendantes et de les comparer à des méthodes classiques de traitement d'image afin de déterminer l'intérêt réel de l'approche proposée.

\vspace{0.3cm}
\noindent
La problématique de ce projet est donc la suivante : \textbf{comment automatiser la segmentation et la quantification de la zone de plaie dans des images de scratch assay à l'aide de techniques de Deep Learning, tout en évaluant leur fiabilité par rapport à une méthode classique de traitement d'image ?}

\vspace{0.3cm}
\noindent
\hspace{1cm}L'objectif principal de ce projet est de développer une solution basée sur l'intelligence artificielle permettant d'automatiser l'analyse des images de scratch assay. Plusieurs objectifs spécifiques ont été définis : préparer et contrôler les données ainsi que les annotations disponibles ; développer un modèle de segmentation basé sur une architecture U-Net avec un encodeur pré-entraîné ; évaluer ses performances sur un ensemble de test indépendant ; comparer les résultats obtenus avec une approche classique basée sur l'analyse de texture et le seuillage d'Otsu ; étudier l'utilisation de plusieurs modèles afin d'améliorer la robustesse de la segmentation et de fournir un indicateur de confiance ; et enfin, développer une application interactive permettant la segmentation, la visualisation, la quantification et la correction manuelle des résultats. Un assistant basé sur une approche de Retrieval-Augmented Generation (RAG) est également intégré afin de faciliter l'interprétation des mesures calculées à partir de la littérature scientifique. Cet assistant intervient uniquement en aval du traitement et ne réalise ni segmentation, ni calcul des mesures, ni diagnostic.

\vspace{0.3cm}
\noindent
\hspace{1cm}Pour conduire ce projet, une démarche structurée basée sur la méthodologie \textbf{CRISP-DM} (\textit{Cross-Industry Standard Process for Data Mining}) a été adoptée. Elle permet d'organiser le travail autour de plusieurs étapes, notamment la compréhension du problème, la compréhension et la préparation des données, la modélisation, l'évaluation et la mise en \oe uvre de la solution. Le projet suit également une démarche empirique et itérative : les différentes décisions techniques sont évaluées à partir de mesures quantitatives sur un jeu de données indépendant lorsque cela est possible. Les résultats obtenus, y compris les limites et les résultats moins satisfaisants, sont pris en compte afin d'éviter de retenir une approche uniquement sur la base d'observations qualitatives.

\vspace{0.3cm}
\noindent
\hspace{1cm}Ce projet a été réalisé au sein du \textbf{Laboratoire de Biologie de la Faculté de Médecine et de Pharmacie de Sousse}, situé à \textbf{Sousse, Tunisie}, dans le cadre d'un stage d'été en intelligence artificielle et Data Science, sous la supervision de \textbf{[Nom du superviseur]}. Le projet s'inscrit dans un contexte de recherche appliquée à l'analyse d'images biologiques, à l'intersection entre les besoins d'un laboratoire de recherche et les techniques d'intelligence artificielle.

\vspace{0.3cm}
\noindent
\hspace{1cm}Le rapport est organisé en cinq chapitres. Le premier chapitre, consacré à l'état de l'art, présente le contexte biologique du scratch assay ainsi que les principales approches utilisées pour l'analyse et la segmentation des images, avec un intérêt particulier pour le Deep Learning, l'architecture U-Net et les méthodes classiques de traitement d'image, et positionne la solution proposée par rapport aux approches existantes. Le deuxième chapitre, consacré à l'analyse et aux spécifications, présente la problématique, les objectifs, les données utilisées, la méthodologie adoptée, ainsi que l'organisme d'accueil et la planification du stage. Le troisième chapitre décrit la conception et l'implémentation de la solution : données, modèle de segmentation, application Streamlit et assistant RAG. Le quatrième chapitre présente le protocole d'évaluation, les métriques utilisées, la comparaison avec l'approche classique ainsi que les résultats expérimentaux et leur analyse. Enfin, le cinquième chapitre présente une synthèse du travail réalisé, les principales limites rencontrées et les perspectives d'amélioration.

% ============================================================
% À partir d'ici : style des chapitres NUMÉROTÉS (Chapitre 1, 2, ...)
% "Chapitre N" + ligne, puis le titre dans un cadre violet — différent
% du style simple (centré, sans cadre) utilisé pour Résumé/Abstract/
% Introduction générale ci-dessus.
% ============================================================
\titleformat{\chapter}[display]
  {\normalfont}
  {\hspace{0.4cm}\bfseries\Large\chaptertitlename~\thechapter\hspace{0.5em}\hrulefill}
  {14pt}
  {\begin{tcolorbox}[colframe=accentbar,colback=white,boxrule=0.8pt,arc=0pt,
      left=10pt,right=10pt,top=8pt,bottom=8pt,width=\textwidth,halign=center]
    \Large\bfseries}
  [\end{tcolorbox}]
\titlespacing*{\chapter}{0pt}{0.28\textheight}{40pt}

% ============================================================
% CHAPITRE 1
% ============================================================
\chapter{État de l'art / Revue bibliographique}

\section{Introduction}

Ce chapitre situe le projet par rapport à la littérature existante. Il présente d'abord le scratch assay lui-même en tant que technique biologique (\ref{sec:scratch}), puis les concepts fondamentaux d'intelligence artificielle et de vision par ordinateur mobilisés (\ref{sec:ia}). Les méthodes de traitement et de segmentation d'image sont ensuite abordées, des approches classiques — seuillage d'Otsu, analyse de texture — aux méthodes par deep learning, avec une attention particulière portée à l'architecture U-Net, au transfer learning et à l'utilisation de ResNet34 comme encodeur, ainsi qu'aux principales métriques utilisées pour évaluer une segmentation. Une revue des travaux appliquant ces méthodes au scratch assay et à des problèmes biologiques proches est ensuite proposée, complétée par les approches méthodologiques non centrales à la segmentation mais mobilisées par le projet — ensemble de modèles et RAG. Le chapitre se conclut par la justification des choix techniques retenus et une synthèse positionnant la solution proposée par rapport à l'existant, en préparation du Chapitre 2 qui détaille les spécifications du projet lui-même. Conformément à l'usage, les résultats chiffrés propres à ce projet ne sont pas présentés ici : ce chapitre décrit ce que propose la littérature, le Chapitre 4 confronte ces repères aux résultats réellement obtenus.

\section{Le scratch assay : principe, protocole et limites}
\label{sec:scratch}

\subsection{Principe biologique}

Le \textit{scratch assay} (ou \textit{wound healing assay}, test de cicatrisation) est un test in vitro qui consiste à créer artificiellement, dans un tapis confluent de cellules adhérentes, une zone linéaire dépourvue de cellules (la « rayure »), puis à observer sa recolonisation progressive au cours du temps \cite{ref13,ref14}. La fermeture résulte principalement de la migration collective des cellules bordant la zone, avec une contribution possible de la prolifération cellulaire selon le type cellulaire et la fenêtre d'observation \cite{ref14}. Le test est utilisé aussi bien en recherche fondamentale sur les mécanismes de migration qu'en criblage de molécules à visée thérapeutique.

\subsection{Protocole expérimental et acquisition des images}

La rayure est généralement créée à l'aide d'une pointe de micropipette raclée sur le tapis cellulaire, puis la zone est imagée par microscopie à différents instants (typiquement toutes les quelques heures) jusqu'à fermeture complète ou quasi-complète \cite{ref13,ref14}. Chaque image constitue un instantané de la progression de la fermeture, ce qui permet en principe une analyse aussi bien ponctuelle (un instant donné) que temporelle (cinétique de fermeture sur plusieurs instants).

\subsection{Méthodes de quantification de la fermeture}

La quantification repose le plus souvent sur la mesure de la surface non colonisée à chaque instant, rapportée à la surface initiale pour obtenir un pourcentage de fermeture \cite{ref9,ref10}, selon une relation de la forme :
\[
\text{Fermeture (\%)} = \frac{S_0 - S_t}{S_0} \times 100
\]
où $S_0$ représente la surface initiale de la plaie et $S_t$ sa surface à l'instant $t$. Certains travaux vont plus loin en estimant une vitesse de migration par intervalle de temps, ou en isolant la contribution de la migration de celle de la prolifération par un suivi de particules ou de trajectoires cellulaires individuelles plutôt qu'une seule mesure de surface globale \cite{ref11}.

\subsection{Sources de variabilité et limites expérimentales}

Le scratch assay est reconnu pour sa simplicité et son faible coût, mais aussi pour plusieurs limites bien documentées \cite{ref13,ref18} : la largeur et la régularité de la rayure dépendent fortement du geste de l'opérateur, ce qui introduit une variabilité inter-opérateur difficile à contrôler ; le test ne mesure que la migration collective planaire sur une surface plane et rigide, sans reproduire les composantes inflammatoires, vasculaires ou matricielles d'une cicatrisation in vivo ; enfin, la qualité de l'image (éclairage, contraste, mise au point) affecte directement la fiabilité de toute mesure ultérieure, qu'elle soit manuelle ou automatisée.

\subsection{Approches de standardisation}

Pour limiter la variabilité opérateur-dépendante, plusieurs travaux proposent de standardiser la création de la rayure elle-même plutôt que seulement son analyse : dispositifs mécaniques ou robotisés reproduisant un geste identique d'un puits à l'autre, gabarits ou peignes multi-pointes, protocoles normalisés d'acquisition d'image \cite{ref18}. Cette piste est complémentaire — et non concurrente — de l'automatisation de l'analyse d'image visée par ce projet : elle agit en amont (la qualité de la rayure produite), là où ce projet agit en aval (la mesure à partir d'une image donnée, quelle que soit sa provenance).

\section{Concepts fondamentaux de l'intelligence artificielle}
\label{sec:ia}

\subsection{Intelligence artificielle et apprentissage automatique}

L'intelligence artificielle désigne, au sens large, la conception de systèmes capables de réaliser des tâches associées à l'intelligence humaine. L'apprentissage automatique (\textit{machine learning}) en est un sous-domaine dans lequel le système n'est pas programmé explicitement pour la tâche, mais apprend un modèle à partir de données.

\subsection{Apprentissage supervisé}

Dans l'apprentissage supervisé, le modèle est entraîné à partir de paires (entrée, sortie attendue) fournies par des exemples annotés — ici, des images de microscopie associées à leur masque de segmentation de référence. C'est le cadre dans lequel s'inscrit la tâche centrale de ce projet : une segmentation binaire (pixel « plaie » / pixel « tapis cellulaire ») apprise à partir d'exemples annotés manuellement, et non une classification d'image entière — la prédiction par pixel permet de calculer directement une surface.

\subsection{Apprentissage non supervisé}

À l'inverse, l'apprentissage non supervisé cherche à extraire une structure de données non annotées (par exemple un regroupement par similarité). Il n'est pas utilisé pour la tâche de segmentation elle-même dans ce projet, mais un principe apparenté — la similarité entre représentations vectorielles (\textit{embeddings}) sans étiquette associée — sous-tend la recherche documentaire par similarité sémantique utilisée par l'assistant RAG (voir \ref{sec:rag}).

\subsection{Deep Learning}

Le Deep Learning est une famille de méthodes d'apprentissage automatique fondées sur des réseaux de neurones à plusieurs couches, capables d'apprendre directement une représentation utile des données brutes (ici, des pixels) sans extraction manuelle de descripteurs — par opposition aux méthodes classiques de traitement d'image, qui reposent sur des descripteurs conçus a priori (seuils, textures).

\subsection{Vision par ordinateur}

La vision par ordinateur regroupe les tâches visant à extraire une information utile d'une image ou d'une vidéo : classification, détection d'objets, et la tâche mobilisée ici, la segmentation — l'attribution d'une catégorie à chaque pixel de l'image.

\section{Traitement et segmentation classique des images}

\subsection{Prétraitement des images}

Avant toute segmentation, les images de microscopie peuvent nécessiter un prétraitement : réduction du bruit, correction du vignettage ou des inhomogénéités d'éclairage (\textit{flat-field}), amélioration locale du contraste (par exemple par égalisation d'histogramme adaptative, CLAHE). Ces étapes visent à réduire la sensibilité des méthodes d'analyse ultérieures aux conditions d'acquisition.

\subsection{Segmentation d'images}

La segmentation consiste à partitionner une image en régions correspondant à des objets ou des zones d'intérêt. Selon la finesse recherchée, on distingue notamment la segmentation par régions (regroupement de pixels similaires) et la segmentation par contours (détection des frontières entre régions).

\subsection{Segmentation sémantique}

La segmentation sémantique attribue une classe à chaque pixel de l'image (ici, deux classes : « plaie » et « tapis cellulaire »), par opposition à la segmentation d'instances qui distinguerait en plus des occurrences individuelles d'un même objet — une distinction non pertinente ici puisque la zone de plaie forme une seule région connexe.

\subsection{Méthodes classiques de segmentation}

Avant l'essor du deep learning, la segmentation reposait sur des règles ou des descripteurs conçus manuellement : seuillage global ou adaptatif de l'intensité, filtrage morphologique, croissance de région, contours actifs, ou transformées par ensembles de niveaux (\textit{level-set}) pour détecter un contour même en conditions de faible contraste.

\subsection{Seuillage d'Otsu et analyse de texture}

La méthode d'Otsu \cite{ref21} détermine automatiquement un seuil d'intensité qui maximise la séparation entre deux classes de pixels, en minimisant la variance intra-classe. Elle suppose implicitement que les deux classes se distinguent par leur intensité — une hypothèse qui ne tient pas toujours pour le scratch assay, où la zone de plaie se distingue du tapis cellulaire non par sa couleur mais par l'absence de texture cellulaire (surface lisse contre surface grainée). Le seuillage d'Otsu peut alors être appliqué non pas directement à l'image, mais à une carte de texture dérivée (par exemple l'écart-type local de l'intensité, ou des descripteurs de texture de type matrice de co-occurrence, GLCM) — une combinaison largement documentée dans les outils historiques d'analyse du scratch assay (voir \ref{sec:travaux}).

\section{Segmentation d'images par Deep Learning}
\label{sec:deep_learning}

La segmentation d'images constitue une tâche importante en vision par ordinateur, particulièrement dans le domaine de l'imagerie biomédicale. Elle consiste à attribuer une classe à chaque pixel d'une image afin d'identifier et de délimiter les régions d'intérêt. Les méthodes de Deep Learning ont progressivement amélioré les performances de segmentation grâce à leur capacité à apprendre automatiquement des représentations adaptées aux données. Dans le cadre de ce projet, une architecture U-Net associée à un encodeur ResNet34 pré-entraîné est utilisée pour segmenter la zone de plaie dans les images issues du scratch assay.

\subsection{Réseaux de neurones convolutifs}

Les réseaux de neurones convolutifs (\textit{Convolutional Neural Networks}, CNN) sont une famille de réseaux de neurones particulièrement adaptée au traitement des images. Contrairement aux méthodes reposant sur des descripteurs conçus manuellement, les CNN apprennent automatiquement des caractéristiques à partir des données d'entrée.

Une couche de convolution applique plusieurs filtres sur l'image afin d'extraire progressivement différentes caractéristiques visuelles. Dans les premières couches, les filtres peuvent apprendre des éléments simples tels que les contours et les textures. Les couches plus profondes permettent ensuite de construire des représentations plus complexes à partir de ces caractéristiques élémentaires.

Cette capacité d'apprentissage hiérarchique rend les CNN particulièrement adaptés aux tâches de vision par ordinateur, notamment la classification, la détection et la segmentation d'images. Dans le cas de la segmentation, l'objectif n'est toutefois pas uniquement de déterminer la classe d'une image entière, mais de produire une prédiction pour chaque pixel.

\subsection{Architecture U-Net}

U-Net~\cite{ref1} est une architecture de réseau de neurones initialement conçue pour la segmentation d'images biomédicales. Son architecture repose sur deux chemins complémentaires : un chemin contractant permettant d'extraire les caractéristiques de l'image et un chemin expansif permettant de reconstruire une carte de segmentation à une résolution proche de celle de l'image originale.

Une caractéristique essentielle de U-Net est l'utilisation de connexions de saut (\textit{skip connections}) entre les niveaux correspondants des deux chemins. Ces connexions permettent de transmettre au décodeur des informations spatiales provenant de l'encodeur, notamment les détails fins nécessaires à la localisation précise des frontières des régions segmentées.

Cette architecture est particulièrement intéressante pour les applications biomédicales, où la délimitation précise des structures d'intérêt est importante et où la quantité de données annotées peut être limitée. Dans ce projet, U-Net est utilisé pour distinguer les pixels correspondant à la zone de plaie des pixels appartenant au tapis cellulaire.

\subsection{Architecture encodeur-décodeur}

L'architecture U-Net peut être considérée comme une architecture de type encodeur-décodeur. L'encodeur constitue le chemin contractant du réseau. Il réduit progressivement la résolution spatiale de l'image tout en augmentant la richesse des caractéristiques extraites. Cette étape permet au réseau de construire une représentation de plus en plus abstraite du contenu de l'image.

Le décodeur constitue le chemin expansif. Il augmente progressivement la résolution des représentations afin de produire une segmentation pixel par pixel. Les connexions de saut entre l'encodeur et le décodeur permettent de combiner les informations sémantiques extraites à basse résolution avec les informations spatiales conservées à des résolutions plus élevées.

Dans le contexte du scratch assay, cette combinaison est importante car la segmentation doit identifier la zone de plaie tout en conservant une délimitation suffisamment précise de ses contours.

\subsection{Transfer Learning}

Le \textit{Transfer Learning}, ou apprentissage par transfert, consiste à réutiliser un modèle ayant déjà appris des représentations sur une base de données source disposant de beaucoup de données (typiquement une classification d'images naturelles sur ImageNet~\cite{ref23}) afin de l'adapter à une nouvelle tâche ou à un nouveau domaine disposant de peu de données. Cette approche est largement documentée en imagerie biomédicale~\cite{ref4,ref5} comme un moyen d'obtenir de bonnes performances malgré un jeu d'entraînement réduit — la littérature met toutefois en garde contre les jeux de données biomédicaux « imparfaits » (petits, bruités, annotés de façon hétérogène), qui restent une contrainte à part entière même avec un encodeur pré-entraîné~\cite{ref5}.

Dans ce projet, cette approche permet de bénéficier de représentations déjà apprises tout en adaptant le réseau aux caractéristiques particulières des images de microscopie du scratch assay, plutôt que d'entraîner l'encodeur entièrement à partir de zéro.

\subsection{ResNet34 comme encodeur}

ResNet34~\cite{ref22} est une architecture de réseau neuronal convolutionnel appartenant à la famille des réseaux résiduels (\textit{Residual Networks}). Elle utilise des connexions résiduelles permettant de transmettre directement certaines informations à travers plusieurs couches du réseau. Cette conception facilite l'entraînement de réseaux relativement profonds et contribue à limiter les problèmes liés à la dégradation des performances lorsque la profondeur du réseau augmente.

Utiliser un ResNet34 pré-entraîné comme encodeur d'un U-Net, plutôt qu'un encodeur « maison » entraîné \textit{from scratch}, est une combinaison documentée dans plusieurs travaux de segmentation biomédicale~\cite{ref3,ref4}. Dans l'architecture utilisée dans ce projet, ResNet34 est employé comme encodeur du U-Net : il est initialisé à partir de poids pré-entraînés, puis adapté aux images de microscopie étudiées. L'encodeur extrait progressivement des caractéristiques visuelles qui sont ensuite transmises au décodeur afin de reconstruire la segmentation.

Le choix de ResNet34 constitue ainsi un compromis entre la capacité de représentation du modèle et sa complexité, ce qui est pertinent dans le contexte d'un jeu de données de taille limitée.

\subsection{Fonctions de perte pour la segmentation}

Lors de l'entraînement d'un modèle de segmentation, une fonction de perte permet de mesurer l'écart entre le masque prédit par le réseau et le masque de référence. Le modèle ajuste ensuite ses paramètres afin de minimiser cette erreur.

Dans le cas d'une segmentation binaire, l'entropie croisée binaire (\textit{Binary Cross-Entropy}, BCE) constitue une fonction de perte couramment utilisée. Calculée indépendamment pour chaque pixel, elle mesure l'écart entre les probabilités prédites et les classes de référence, mais traite les pixels selon leur contribution individuelle à la classification, ce qui peut la rendre moins adaptée aux situations de fort déséquilibre entre les classes — une situation fréquente en segmentation biomédicale, où l'objet d'intérêt (ici, la fine bande de plaie) occupe une faible proportion de l'image.

Pour améliorer la prise en compte du chevauchement entre les régions prédites et les régions de référence, une perte basée sur le coefficient de Dice (voir section~\ref{sec:metriques}) peut également être utilisée. Ce coefficient s'écrit de façon équivalente :
\[
\text{Dice} = \frac{2\,TP}{2\,TP + FP + FN}
\]
où $TP$ représente les vrais positifs, $FP$ les faux positifs et $FN$ les faux négatifs. La \textit{Dice Loss}, dérivée directement de ce coefficient, pénalise le manque de recouvrement global entre prédiction et référence plutôt que l'erreur pixel par pixel.

Dans ce projet, les deux composantes sont combinées afin de former une fonction de perte de type Dice + BCE, ce qui permet de bénéficier à la fois d'un signal d'apprentissage stable au niveau des pixels et d'une sensibilité au chevauchement global entre la zone de plaie prédite et le masque de référence.

\section{Évaluation des méthodes de segmentation}
\label{sec:metriques}

La littérature met en garde contre une évaluation naïve de la segmentation, qui choisirait des métriques par habitude plutôt qu'en fonction du problème posé \cite{ref6,ref7}. Les métriques suivantes sont les plus couramment utilisées pour la segmentation binaire et sont, pour cette raison, retenues pour ce projet (leur valeur mesurée sur ce projet est présentée au Chapitre 4, pas ici).

\subsection{Dice Score}

Le coefficient de Dice mesure le recouvrement entre le masque prédit $P$ et le masque de référence $G$ : $\text{Dice} = \dfrac{2 |P \cap G|}{|P| + |G|}$. Il varie de 0 (aucun recouvrement) à 1 (recouvrement parfait) et pondère implicitement davantage les faux positifs/négatifs sur les petits objets que sur les grands.

\subsection{Intersection over Union}

L'IoU (ou indice de Jaccard) mesure le rapport entre l'intersection et l'union des deux masques : $\text{IoU} = \dfrac{|P \cap G|}{|P \cup G|}$. Il est mathématiquement lié au Dice ($\text{IoU} = \text{Dice} / (2 - \text{Dice})$) mais pénalise proportionnellement plus fortement les erreurs que le Dice.

\subsection{Précision et rappel}

La précision ($|P \cap G| / |P|$) mesure la proportion de pixels prédits comme plaie qui le sont réellement ; le rappel ($|P \cap G| / |G|$) mesure la proportion de pixels réellement « plaie » effectivement détectés. Un modèle peut obtenir un bon Dice tout en étant déséquilibré entre les deux (par exemple, sur-segmenter légèrement pour maximiser le rappel au prix de la précision) — d'où l'intérêt de les rapporter séparément plutôt que de se limiter au Dice.

\subsection{Distance de Hausdorff}

Contrairement aux métriques précédentes, fondées sur le recouvrement de surface, la distance de Hausdorff mesure un écart de contour : la plus grande distance, parmi tous les points d'un contour, à son point le plus proche sur l'autre contour. C'est une métrique de pire cas, très sensible à un seul point de contour mal placé, même si le recouvrement global est excellent par ailleurs \cite{ref7}. L'article introduisant la mesure \textit{clDice} \cite{ref8} propose une alternative topologique spécifiquement conçue pour les structures fines et allongées (initialement les réseaux vasculaires), une géométrie proche de celle d'une bande de plaie fine et longue.

\subsection{Erreur de surface}

Spécifique au cas d'usage biologique du scratch assay, l'erreur de surface relative exprime directement l'écart de mesure dans l'unité utile au biologiste (une surface, convertible en pourcentage de fermeture), plutôt qu'un score de recouvrement abstrait :
\[
\text{Erreur de surface (\%)} = \frac{|S_{\text{prédite}} - S_{\text{référence}}|}{S_{\text{référence}}} \times 100
\]
où $S_{\text{prédite}}$ représente la surface calculée à partir du masque prédit et $S_{\text{référence}}$ la surface obtenue à partir du masque de référence. Plusieurs travaux de quantification du scratch assay rapportent directement ce type d'erreur plutôt que le seul Dice \cite{ref9,ref10}.

\section{Travaux similaires et état de la recherche}
\label{sec:travaux}

\subsection{Automatisation de l'analyse du scratch assay}

Plusieurs outils dédiés à l'analyse d'images du scratch assay sans recours au deep learning ont été développés et restent utilisés en laboratoire. Parmi eux figurent le plugin ImageJ \textit{Wound Healing Size Tool}, l'outil \textit{TScratch}, fondé sur une transformée en curvelets \cite{ref15}, ainsi que \textit{Fiji} \cite{ref16}, qui constitue une plateforme généraliste d'analyse d'images biologiques. Des outils plus spécifiquement automatisés, comme \textit{AIM} \cite{ref17}, combinent le seuillage adaptatif et des opérations de traitement morphologique afin de produire une mesure de fermeture sans intervention manuelle répétée.

Ces outils restent appréciés pour leur simplicité et leur faible coût computationnel, mais nécessitent généralement un réglage manuel de certains paramètres et restent sensibles aux variations d'éclairage et de contraste.

\subsection{Segmentation biologique par Deep Learning}

Au-delà du scratch assay, le U-Net et ses variantes, notamment U-Net++, ainsi que les architectures intégrant des encodeurs pré-entraînés comme TernausNet \cite{ref4}, sont largement utilisés pour la segmentation d'images biomédicales. Ces approches sont particulièrement intéressantes dans des contextes où les données annotées disponibles sont limitées \cite{ref5}, une contrainte également rencontrée dans de nombreux problèmes d'imagerie biologique.

Le deep learning permet d'apprendre automatiquement des caractéristiques pertinentes à partir d'images annotées, mais les performances obtenues restent liées à la qualité et à la représentativité des données d'apprentissage ainsi qu'à la capacité du modèle à généraliser à de nouvelles images.

\subsection{Utilisation de U-Net dans l'analyse du scratch assay}

Deux travaux directement pertinents pour ce projet ont été identifiés. Le premier applique un U-Net à l'analyse du scratch assay sur un jeu de données de plus de 400 images provenant d'un même laboratoire et rapporte un coefficient de Dice compris entre 0,958 et 0,968 \cite{ref2}. Le second applique une architecture comparable, fondée sur un U-Net avec un encodeur pré-entraîné, à des micrographies de cellules souches mésenchymateuses en contraste de phase. Bien que cette application ne porte pas directement sur le scratch assay, cette modalité d'imagerie présente certaines caractéristiques proches. Le travail rapporte un Dice de 0,876 \cite{ref3}.

Un troisième travail combine une segmentation par U-Net++ avec différentes méthodes de modélisation temporelle de la trajectoire de fermeture, notamment les forêts aléatoires, ARIMA et un réseau convolutif temporel \cite{ref12}. Par ailleurs, un suivi de particules individuelles a également été proposé afin d'analyser la cinétique de fermeture image par image \cite{ref11}.

\subsection{Comparaison des approches existantes}

Les résultats rapportés par ces travaux dépendent de plusieurs facteurs, notamment de la taille et de l'homogénéité du jeu de données d'entraînement, du protocole expérimental et de la modalité d'imagerie. Le travail disposant de plus de 400 images issues d'un protocole unique \cite{ref2} rapporte un Dice sensiblement supérieur à celui obtenu dans l'étude portant sur un domaine d'imagerie proche avec un jeu de données plus restreint \cite{ref3}. Toutefois, cet écart ne peut pas être attribué uniquement à la taille du jeu de données, puisque les deux études diffèrent également par leurs conditions expérimentales et leurs modalités d'imagerie.

Ces deux références constituent néanmoins des repères pertinents pour la comparaison des performances et seront utilisées au Chapitre 4 afin de situer les résultats obtenus dans le cadre de ce projet.

\subsection{Limites des travaux existants}

Peu de travaux identifiés couplent une segmentation automatique à un suivi temporel multipoint de la fermeture, par opposition à une mesure ponctuelle ou à une comparaison entre deux instants \cite{ref11,ref12}. Par ailleurs, le travail directement consacré au scratch assay \cite{ref2} ne rapporte pas d'évaluation sur des images issues d'un laboratoire ou d'un microscope différent de celui utilisé pour l'entraînement. La question de la généralisation inter-laboratoire reste ainsi ouverte dans ce domaine de recherche.

Ces limites motivent, dans le cadre de ce projet, l'utilisation d'un jeu de test indépendant pour l'évaluation, la comparaison avec une méthode de référence classique et l'étude d'approches complémentaires, notamment l'ensemble de modèles et un indicateur d'incertitude, présentées dans la section suivante.

\section{Approches complémentaires}
\label{sec:rag}

En complément des méthodes de segmentation présentées précédemment, certaines approches peuvent être mobilisées afin d'améliorer la robustesse des prédictions et de faciliter l'interprétation des résultats. Dans ce projet, elles interviennent en complément du modèle principal de segmentation et concernent l'utilisation d'un ensemble de modèles, un indicateur d'incertitude ainsi que les approches fondées sur le Retrieval-Augmented Generation (RAG) et les modèles de langage.

\subsection{Ensemble de modèles et amélioration de la robustesse}

Combiner plusieurs modèles indépendants constitue une technique classique d'ensemble en apprentissage automatique. Dans le cas de la segmentation d'images, les prédictions de plusieurs modèles peuvent notamment être combinées par vote majoritaire pixel par pixel. Cette stratégie permet de réduire la dépendance à un modèle unique et peut améliorer la robustesse des prédictions lorsque les modèles présentent des erreurs différentes.

La diversité entre les modèles peut notamment provenir de variations dans leur initialisation ou dans les conditions d'apprentissage. Lorsque leurs erreurs ne sont pas parfaitement corrélées, leur combinaison peut conduire à une prédiction plus stable que celle d'un modèle individuel.

Dans le cadre de ce projet, plusieurs modèles entraînés avec des initialisations différentes sont réutilisés afin de construire une prédiction par ensemble. Les résultats expérimentaux associés à cette approche sont présentés au Chapitre 4.

\subsection{Indicateur d'incertitude basé sur l'accord entre modèles}

Au-delà de l'amélioration de la prédiction, le désaccord entre les modèles d'un ensemble peut être exploité comme un indicateur d'incertitude par pixel. Lorsqu'une majorité de modèles produit la même prédiction, l'accord est élevé. À l'inverse, une divergence entre les modèles indique une région pour laquelle la prédiction est moins stable.

Cette idée est utilisée dans différentes approches de quantification de l'incertitude en imagerie médicale fondées sur le Deep Learning \cite{ref19,ref20}. Elle permet notamment de compléter la prédiction binaire par une information indiquant les régions où le modèle présente davantage d'incertitude.

Dans ce projet, le désaccord entre les modèles de l'ensemble est utilisé pour produire une carte mettant en évidence les zones incertaines. La relation entre ce signal de désaccord et les performances de segmentation est étudiée expérimentalement au Chapitre 4.

\subsection{Retrieval-Augmented Generation (RAG)}

Le Retrieval-Augmented Generation (RAG) \cite{ref24} est une approche qui combine un modèle de langage avec un mécanisme de recherche dans une base documentaire. Les documents pertinents sont récupérés en fonction de la question posée, notamment à l'aide de représentations vectorielles (\textit{embeddings}), puis fournis au modèle de langage afin de générer une réponse contextualisée.

Cette approche permet d'ancrer les réponses dans un corpus documentaire sélectionné plutôt que de dépendre uniquement des connaissances internes du modèle de langage. Elle est ainsi particulièrement intéressante dans des domaines scientifiques où les réponses doivent pouvoir être reliées à des sources documentaires.

Dans le cadre de ce projet, le RAG est utilisé avec une base documentaire locale regroupant des articles scientifiques relatifs à la cicatrisation et au scratch assay. Il constitue une couche complémentaire destinée à fournir au modèle de langage des informations issues de cette documentation.

\subsection{Utilisation des LLM pour l'interprétation scientifique}

Les grands modèles de langage (\textit{Large Language Models}, LLM) peuvent être utilisés pour reformuler des résultats quantitatifs, répondre à des questions en langage naturel et mettre en relation des observations avec des informations issues de la littérature scientifique. Cependant, leur utilisation dans un contexte scientifique nécessite de prendre en compte le risque de génération d'informations incorrectes mais plausibles, souvent désigné sous le terme d'\textit{hallucination}.

Le RAG peut contribuer à limiter ce risque en fournissant au modèle des documents pertinents servant de contexte, mais il ne garantit pas à lui seul l'exactitude de toutes les réponses générées. Une séparation claire entre le calcul des mesures et leur interprétation constitue donc une précaution importante.

Dans ce projet, les mesures quantitatives sont calculées en amont par les méthodes de traitement et d'analyse implémentées en Python. Le LLM intervient ensuite pour faciliter l'interprétation textuelle de ces résultats, en s'appuyant sur les documents récupérés par le système RAG. Les détails de cette architecture et les mécanismes de contrôle mis en place sont présentés au Chapitre 3.

\section{Justification de l'approche adoptée}

Les éléments présentés dans les sections précédentes permettent de justifier les choix méthodologiques retenus pour répondre à la problématique du projet. L'objectif n'est pas uniquement d'automatiser la mesure de la fermeture de la plaie, mais de construire une approche capable de produire des segmentations reproductibles, d'évaluer leur fiabilité et de faciliter l'interprétation des résultats dans un contexte de recherche en biologie.

\subsection{Choix du Deep Learning}

La revue ci-dessus (1.4-1.5) montre que les méthodes classiques de segmentation, bien qu'utiles et peu coûteuses, reposent sur des descripteurs conçus a priori (intensité, texture) dont la validité dépend fortement des conditions d'acquisition — une limite directement pertinente pour un jeu d'images de microscopie aux conditions d'éclairage variables. Le Deep Learning, en apprenant directement une représentation à partir des données, est une réponse documentée à cette limite dans le champ plus large de la segmentation biomédicale (1.7.2).

\subsection{Choix de U-Net et de ResNet34}

Le U-Net constitue une architecture largement utilisée pour la segmentation biomédicale, notamment lorsque le jeu de données annotées est de taille réduite (1.5.2), et les travaux portant directement sur le scratch assay ou sur des modalités proches (1.7.3) constituent des éléments en faveur de sa pertinence pour ce domaine précis. Un encodeur ResNet34 pré-entraîné sur ImageNet, plutôt qu'un encodeur entraîné \textit{from scratch}, est cohérent avec la pratique documentée en transfer learning biomédical (1.5.4, 1.5.5) — une pratique particulièrement adaptée à un jeu de données d'entraînement de taille modeste (voir Chapitre 2, jeu de données).

\subsection{Intérêt de la comparaison avec une méthode classique}

Plutôt que de supposer a priori la supériorité du Deep Learning, conserver une baseline classique (seuillage sur une carte de texture, dans l'esprit des outils historiques du domaine, 1.7.1) permet de mesurer objectivement l'apport réel du Deep Learning pour ce projet précis, plutôt que de s'appuyer uniquement sur des résultats publiés dans d'autres conditions (voir Chapitre 4).

\subsection{Intérêt de l'ensemble de modèles}

L'ensemble de modèles est motivé par la volonté de réduire la dépendance à une seule initialisation et par l'intérêt d'un indicateur de désaccord par pixel au-delà de la seule prédiction binaire (1.8.2). Cette information est utile dans un outil destiné à assister un biologiste, qui peut ainsi disposer d'une indication sur les régions où les prédictions des modèles sont moins concordantes, plutôt que de considérer la prédiction automatique comme infaillible.

\subsection{Intérêt du RAG pour l'interprétation scientifique}

Le choix d'un RAG plutôt que d'un LLM seul pour l'interprétation des résultats découle directement du risque de fabrication d'information documenté en 1.8.4 : ancrer la réponse sur des extraits réellement récupérés est une réponse partielle mais nécessaire à ce risque, ce qui a motivé le développement d'une méthodologie de test spécifique pour cette partie du projet (voir Chapitre 3).

\subsection{Synthèse : solution proposée}

Le tableau~\ref{tab:solution_proposee} résume la position de la solution développée par rapport aux approches présentées dans ce chapitre.

\begin{small}
\begin{longtable}{L{3.4cm} L{5.4cm} L{5.2cm}}
\caption{Positionnement de la solution proposée par rapport aux approches existantes.}
\label{tab:solution_proposee} \\
\toprule
\textbf{Approche} & \textbf{Principe} & \textbf{Limites relevées} \\
\midrule
\endfirsthead
\toprule
\textbf{Approche} & \textbf{Principe} & \textbf{Limites relevées} \\
\midrule
\endhead
Mesure manuelle (ImageJ) & Tracé du contour et mesure de la surface image par image & Lente sur de nombreuses images, dépendante de l'opérateur \\
Outils classiques (Wound Healing Size Tool, TScratch, AIM) \cite{ref15,ref17} & Seuillage, texture ou curvelets, opérations morphologiques & Réglage manuel de paramètres, sensibilité à l'éclairage et au contraste \\
U-Net appliqué au scratch assay \cite{ref2} & Segmentation par deep learning sur plus de 400 images & Pas d'évaluation rapportée sur un autre laboratoire ou microscope \\
\textbf{Solution proposée} & U-Net + ResNet34 pré-entraîné, comparé à une baseline classique ; carte de confiance par ensemble de modèles ; suivi temporel de la fermeture ; résumé scientifique assisté par RAG, avec synthèse vocale & Jeu de 97 images, pas de validation externe (voir Chapitre 5) \\
\bottomrule
\end{longtable}
\end{small}

\section{Conclusion}

Ce chapitre a présenté le scratch assay comme technique biologique, les concepts d'intelligence artificielle et de segmentation mobilisés par le projet, ainsi qu'une revue des travaux existants sur l'automatisation de son analyse — qu'elle repose sur des méthodes classiques ou sur le Deep Learning. Cette revue met en évidence à la fois l'intérêt du Deep Learning pour ce type de tâche et les limites encore ouvertes du champ de recherche (généralisation inter-laboratoire, suivi temporel multi-points), qui justifient les choix techniques retenus pour ce projet. Le chapitre suivant précise, à partir de ce cadre, la problématique, les objectifs et les spécifications concrètes du projet.

% ============================================================
% CHAPITRE 2
% ============================================================
\chapter{Analyse et spécifications du projet}

\section{Définition du problème}

Étant donné une série d'images de microscopie d'un même scratch assay prises à différents instants, le système doit produire automatiquement :
\begin{enumerate}
\item le masque binaire de la zone non colonisée (la « plaie ») pour chaque image ;
\item la surface correspondante (en pixels\textsuperscript{2}, convertible en cm\textsuperscript{2} si un ratio de calibration est fourni) ;
\item le pourcentage de fermeture par rapport à l'instant initial, ainsi que son évolution au fil des différents instants observés ;
\item un résumé scientifique en langage naturel de cette évolution, replacé dans le contexte de la littérature scientifique pertinente.
\end{enumerate}

\section{Objectifs fonctionnels}

\begin{itemize}
\item Segmenter une image de scratch assay et en extraire la surface non colonisée.
\item Permettre une correction manuelle du masque (suppression de fragments, correction au pinceau) lorsque la prédiction automatique est imparfaite.
\item Suivre plusieurs images d'une même série temporelle et calculer l'évolution de la fermeture.
\item Comparer un masque prédit à un masque de référence (métriques de qualité).
\item Générer un résumé scientifique de l'évolution observée, appuyé sur une base documentaire.
\item Permettre au biologiste d'écouter le résumé scientifique (synthèse vocale), par exemple pendant une manipulation au laboratoire, sans avoir à le relire à l'écran.
\item Fournir un indicateur d'incertitude basé sur l'accord entre modèles (option d'ensemble de modèles), pour signaler à l'utilisateur les zones de désaccord entre modèles plutôt que de présenter un résultat unique sans indication de fiabilité.
\end{itemize}

\section{Objectifs non fonctionnels}

\begin{itemize}
\item \textbf{Reproductibilité} : split d'entraînement/validation/test fixé par graine aléatoire, expériences tracées (MLflow), résultats de robustesse inter-seed mesurés plutôt que supposés.
\item \textbf{Fonctionnement local / sans dépendance externe payante} : le LLM tourne en local via Ollama, sans appel à une API propriétaire facturée à l'usage — contrainte pertinente pour un usage académique ou en laboratoire sans budget API dédié.
\item \textbf{Fiabilité de la génération de texte} : aucune valeur numérique présentée au biologiste ne doit être calculée ou modifiée par le LLM — un objectif non fonctionnel central du projet, traité comme une contrainte d'architecture (voir Chapitre 3) plutôt que comme une simple bonne pratique de prompt.
\item \textbf{Utilisabilité} : interface graphique (Streamlit) ne nécessitant aucune compétence en programmation de l'utilisateur final (biologiste).
\item \textbf{Portée clairement délimitée} : outil de quantification pour la recherche in vitro, explicitement non destiné à un usage diagnostique ou clinique (voir Chapitre 5, limites).
\end{itemize}

\section{Jeu de données}

\begin{itemize}
\item \textbf{Source :} 97 images de microscopie de scratch assay, fournies par le laboratoire d'accueil~\cite{ref41}.
\item \textbf{Annotation :} chaque image est annotée manuellement par un contour tracé à la main directement sur la photographie ; ces annotations ne constituent pas directement des masques binaires exploitables par le modèle.
\item \textbf{Préparation :} un script dédié reconstruit les masques binaires à partir de ces contours tracés à la main : détection du trait par couleur/teinte, pontage des discontinuités du trait (effacé par endroits à cause de la compression JPEG), fermeture morphologique sur les bords d'image, puis remplissage par \textit{flood-fill} de la zone délimitée.
\item \textbf{Nettoyage :} une image est exclue du jeu de données final, son tracé manuel étant incomplet même après ce traitement.
\item \textbf{Split :} séparation reproductible (graine fixée) en 68 images d'entraînement, 14 de validation, 14 de test. Le jeu de test est tenu à l'écart aussi bien de l'entraînement que de la sélection du meilleur modèle (contrairement au jeu de validation), afin que l'évaluation finale fournisse une estimation indépendante des performances de généralisation.
\item \textbf{Variables/format :} images RGB de microscopie (résolution native variable selon l'image), masques binaires de même résolution. Pour la fonctionnalité de suivi temporel, chaque image analysée peut être associée à un instant d'observation exprimé en heures, déduit du nom de fichier ou fourni explicitement, ainsi qu'à un identifiant d'image. Ces informations sont utilisées pour organiser le suivi et calculer l'évolution de la fermeture, et ne constituent pas des variables d'entrée du modèle de segmentation.
\end{itemize}

\section{Méthode d'analyse (pipeline global)}
\label{sec:pipeline}

\begin{figure}[h]
\centering
\includegraphics[width=0.7\textwidth]{images/figure_pipeline.png}
\caption{Pipeline global du système.}
\label{fig:pipeline}
\end{figure}

Le pipeline est volontairement stratifié en couches distinctes, chacune ayant une responsabilité clairement définie :

\begin{enumerate}
\item \textbf{Couche vision} : produit automatiquement, à partir de l'image d'entrée, un masque de segmentation évalué par des métriques d'image classiques (Chapitre 4). Une option d'ensemble de modèles permet également d'obtenir un indicateur d'incertitude fondé sur le désaccord entre les prédictions.
\item \textbf{Couche quantification} : transforme ce masque en grandeurs numériques (surface, pourcentage de fermeture, évolution entre les instants), calculées uniquement en code (Python), jamais par le LLM.
\item \textbf{Couche interprétation} : reçoit ces grandeurs déjà calculées en entrée, puis les met en forme et les contextualise, sans effectuer elle-même de calcul numérique ; des mécanismes de contrôle visent à limiter la génération de valeurs non présentes dans les données calculées (voir Chapitre 3.4).
\item \textbf{Couche présentation} : interface Streamlit reliant les trois couches précédentes pour l'utilisateur final.
\end{enumerate}

La figure~\ref{fig:architecture} présente l'architecture générale du système et met en évidence la séparation entre ces différentes couches fonctionnelles.

\begin{figure}[h]
\centering
\includegraphics[width=\textwidth]{images/figure_architecture.png}
\caption{Architecture générale du système proposé. Le diagramme reprend la stratification en quatre couches décrite ci-dessus : la couche vision produit le masque de segmentation (complétée par l'ensemble de modèles, qui fournit un indicateur d'incertitude par désaccord) ; les couches quantification et interprétation transforment puis contextualisent ces résultats sans jamais se substituer l'une à l'autre ; la couche présentation (Streamlit) est le seul point de contact avec l'utilisateur final.}
\label{fig:architecture}
\end{figure}

\section{Méthodologie de conduite du projet : cadre CRISP-DM}

Au-delà du pipeline technique décrit en \ref{sec:pipeline}, la conduite du projet dans son ensemble a suivi le cadre \textbf{CRISP-DM} (\textit{Cross-Industry Standard Process for Data Mining}), une méthodologie standard de conduite de projet de science des données en six phases. Le tableau~\ref{tab:crispdm} fait correspondre chaque phase à son contenu concret dans ce projet et au chapitre qui la détaille :

\begin{longtable}{p{3.2cm} p{8.5cm} p{2.3cm}}
\caption{Correspondance entre les phases CRISP-DM et le contenu du projet.}
\label{tab:crispdm} \\
\toprule
\textbf{Phase CRISP-DM} & \textbf{Contenu du projet} & \textbf{Détaillé en} \\
\midrule
\endfirsthead
\toprule
\textbf{Phase CRISP-DM} & \textbf{Contenu du projet} & \textbf{Détaillé en} \\
\midrule
\endhead
\textbf{1. Compréhension du métier} & Remplacer le tracé manuel du scratch assay, lent et peu reproductible, par une mesure automatique ; évaluer si le deep learning apporte un gain réel par rapport à une méthode classique ; fournir une interprétation scientifique prudente des résultats, sans jamais se substituer au biologiste & Introduction générale ; 2.1--2.3 ; 2.7 \\
\textbf{2. Compréhension des données} & 97 images de microscopie de scratch assay, annotées par un contour tracé à la main (pas un masque directement exploitable) ; 48 articles scientifiques sélectionnés manuellement pour la base documentaire du RAG & 2.4 ; 3.4 \\
\textbf{3. Préparation des données} & Reconstruction des masques binaires à partir des contours tracés (détection de trait, pontage des discontinuités, remplissage) ; exclusion d'une image à l'annotation incomplète ; split reproductible train/val/test (68/14/14) ; augmentation de données à l'entraînement ; curation et traçabilité des décisions d'inclusion/exclusion des articles du corpus documentaire & 2.4 ; 3.3 \\
\textbf{4. Modélisation} & U-Net + encodeur ResNet34 pré-entraîné, entraînement en deux phases, perte Dice+BCE ; baseline classique de comparaison (seuillage de texture) ; ensemble de 4 modèles indépendants ; Test-Time Augmentation ; architecture LLM/RAG à deux appels séparés (mesures / littérature) & 3.1--3.2, 3.4 \\
\textbf{5. Évaluation} & Métriques de recouvrement et d'erreur de surface sur un jeu de test tenu à l'écart ; test de significativité de Wilcoxon et intervalle de confiance bootstrap ; étude de robustesse inter-seed ; analyse d'erreur qualitative et quantitative par catégorie ; benchmark de retrieval documentaire ; tests de régression sur la fiabilité de la génération LLM ; vérification de la synthèse vocale dans le navigateur & Chapitre 4 \\
\textbf{6. Déploiement} & Application Streamlit en deux onglets, avec synthèse vocale du résumé, utilisable sans compétence de programmation ; conteneurisation Docker ; registre de modèles MLflow avec promotion opt-in par comparaison de Dice, sans échec de l'entraînement en cas de problème de registre ; intégration continue (tests + build Docker + test de fumée) --- sans déploiement automatisé réel, portée volontairement limitée & 3.6 ; Chapitre 5 (limites) \\
\bottomrule
\end{longtable}

\textbf{Nature itérative de la démarche.} Contrairement à une application linéaire et unique de ces six phases, plusieurs cycles ont été nécessaires en pratique — en particulier entre modélisation et évaluation (l'étude de robustesse inter-seed et l'ensemble de modèles ont été ajoutés après une première évaluation ayant révélé l'instabilité du Hausdorff, voir Chapitre 4), et entre préparation des données et évaluation côté RAG (l'élargissement ciblé du corpus documentaire, voir 3.4, a été motivé par des lacunes de retrieval constatées lors de l'évaluation). Cette itération répétée entre phases, plutôt qu'un enchaînement strictement séquentiel, est elle-même conforme à l'esprit de CRISP-DM, qui présente ce cycle comme itératif et non linéaire.

\section{Organisme d'accueil et planification du stage}

\textbf{Organisme d'accueil.} Le stage s'est déroulé au sein du \textbf{Laboratoire de Biologie de la Faculté de Médecine et de Pharmacie de Sousse} (Tunisie). Le laboratoire mène des expériences in vitro, dont des scratch assays suivis par microscopie, et a fourni le jeu de 97 images annotées utilisé dans ce projet~\cite{ref41}. Le besoin exprimé était de réduire le temps passé à délimiter et mesurer manuellement la zone de plaie sur chaque image, tout en conservant un contrôle visuel du résultat. % À COMPLÉTER : axes de recherche du laboratoire, équipe, rôle du superviseur.

\textbf{Planification.} La figure~\ref{fig:gantt} présente le déroulement du projet. Les dates sont reconstituées à partir de l'historique des versions du code (git) et des dates de création des fichiers de données et de résultats, et non d'un planning prévisionnel. Plusieurs phases se chevauchent, conformément au caractère itératif de la démarche CRISP-DM décrit ci-dessus : l'évaluation a par exemple conduit à ajouter l'étude multi-graines et l'ensemble de modèles, et le travail sur le RAG s'est poursuivi pendant le développement de l'application.

\begin{figure}[htbp]
\centering
\includegraphics[width=\textwidth]{images/figure_gantt.png}
\caption{Diagramme de Gantt du projet (dates reconstituées à partir de l'historique git et des fichiers).}
\label{fig:gantt}
\end{figure}

\section{Conclusion}

Ce chapitre a permis de définir le problème traité, les objectifs fonctionnels et non fonctionnels du système, ainsi que les données utilisées et leur préparation. L'architecture générale du pipeline a également été présentée, en distinguant les différentes couches consacrées à la segmentation, à la quantification, à l'interprétation et à la présentation des résultats.

La démarche CRISP-DM adoptée pour la conduite du projet a permis de structurer les différentes étapes du développement tout en conservant une approche itérative ; le diagramme de Gantt en montre le déroulement réel, au sein du laboratoire d'accueil. Ces éléments constituent la base nécessaire à la conception et à l'implémentation détaillée du système, présentées dans le chapitre suivant.

\chapter{Conception et implémentation du système}

\section{Architecture du système}

\textbf{Architecture générale.} Le système suit la stratification en quatre couches déjà présentée au Chapitre 2 (vision, quantification, interprétation, présentation ; voir 2.5) : chaque couche correspond ici à un ou plusieurs modules concrets du dépôt.

\begin{figure}[h]
\centering
\includegraphics[width=0.85\textwidth]{images/figure_modules.png}
\caption{Modules du système et leurs interactions.}
\label{fig:modules}
\end{figure}

\textbf{Modules du système.}

\begin{longtable}{p{2.3cm} p{5.5cm} p{5.5cm}}
\caption{Modules du dépôt et fichiers clés associés.}
\label{tab:modules} \\
\toprule
\textbf{Module} & \textbf{Rôle} & \textbf{Fichiers clés} \\
\midrule
\endfirsthead
\toprule
\textbf{Module} & \textbf{Rôle} & \textbf{Fichiers clés} \\
\midrule
\endhead
\texttt{src/} & Préparation des données, modèle, entraînement, évaluation, mesure, cinétique & \texttt{model.py}, \texttt{train.py}, \texttt{dataset.py}, \texttt{evaluate.py}, \texttt{measure.py}, \texttt{kinetics.py}, \texttt{baseline.py}, \texttt{ensemble.py} \\
\texttt{chatbot/} & Assistant IA (RAG) : recherche documentaire, génération de résumé, contrôles après génération & \texttt{chatbot.py}, \texttt{index\_articles.py}, \texttt{evaluate\_retrieval.py} \\
\texttt{app.py} & Interface utilisateur (Streamlit, deux onglets, synthèse vocale du résumé) & --- \\
\texttt{tests/} & Tests automatisés (unitaires + réels contre le LLM) & \texttt{test\_measure.py}, \texttt{test\_metrics.py}, \texttt{test\_ensemble.py}, \texttt{test\_chatbot\_*.py} \\
\bottomrule
\end{longtable}

\textbf{Interactions entre les modules.} \texttt{app.py} orchestre les trois autres modules pour l'utilisateur final : il appelle \texttt{src/} pour la segmentation et la quantification (couches vision et quantification), puis \texttt{chatbot/} pour l'interprétation scientifique (couche interprétation), en ne lui transmettant que des grandeurs déjà calculées (voir 3.4). \texttt{tests/} valide indépendamment \texttt{src/} et \texttt{chatbot/}, y compris par des tests réels contre le LLM local (voir 3.4).

\textbf{Technologies et outils utilisés.} Python, avec les bibliothèques principales suivantes : \texttt{torch} / \texttt{torchvision}~\cite{ref25}, \texttt{segmentation-models-pytorch}~\cite{ref26}, \texttt{albumentations}~\cite{ref27}, \texttt{opencv-python}~\cite{ref28}, \texttt{numpy}~\cite{ref29}, \texttt{scipy}~\cite{ref30} (test de Wilcoxon, distance de Hausdorff), \texttt{matplotlib}~\cite{ref31}, \texttt{streamlit}~\cite{ref32} (+ \texttt{streamlit-drawable-canvas}, \texttt{streamlit-image-coordinates}), \texttt{chromadb}~\cite{ref33} et \texttt{rank\_bm25}~\cite{ref34,ref35} (recherche documentaire), \texttt{ollama}~\cite{ref36} (client LLM local), synthèse vocale native du navigateur (\textit{Web Speech API}~\cite{ref37}, sans bibliothèque supplémentaire), \texttt{mlflow}~\cite{ref38} (suivi d'expériences), \texttt{pytest}~\cite{ref39} (tests). L'entraînement et le suivi d'expériences utilisent MLflow (métriques par époque, checkpoints et logs sauvegardés comme artefacts pour chaque run), avec une intégration continue (tests automatisés) sur GitHub Actions à chaque modification du code (voir 3.6).

\textbf{Configuration de l'environnement.} Le développement, l'entraînement et les évaluations ont été réalisés sur un poste unique, décrit au tableau~\ref{tab:environnement}. L'entraînement utilise automatiquement le GPU lorsqu'il est disponible (\texttt{src/train.py}), sinon le CPU ; l'image Docker, elle, est volontairement limitée au CPU (voir 3.6). Le fichier \texttt{requirements.txt} ne fige pas les versions : celles du tableau sont les versions installées et vérifiées sur ce poste.

\begin{longtable}{L{5.5cm} L{8.5cm}}
\caption{Configuration matérielle et logicielle utilisée.}
\label{tab:environnement} \\
\toprule
\textbf{Élément} & \textbf{Version / caractéristique} \\
\midrule
\endfirsthead
\toprule
\textbf{Élément} & \textbf{Version / caractéristique} \\
\midrule
\endhead
Système d'exploitation & Windows 11 (64 bits) \\
Processeur & Intel Core i5-12450HX (12\textsuperscript{e} génération) \\
Mémoire vive & 16 Go \\
Carte graphique & NVIDIA GeForce RTX 2050 (CUDA 12.4 via PyTorch) \\
Langage & Python 3.11.9 \\
Deep learning & PyTorch 2.6.0, torchvision 0.21.0, segmentation-models-pytorch 0.5.0 \\
Données et images & albumentations 2.0.8, OpenCV 4.13.0, NumPy 2.4.6, Pillow 12.2.0 \\
Statistiques et figures & SciPy 1.16.3, Matplotlib 3.10.9 \\
Recherche documentaire & ChromaDB 1.5.9, rank\_bm25 0.2.2 \\
Modèle de langage local & Ollama 0.34.4 (serveur), client Python ollama 0.6.2, modèle Mistral 7B \\
Interface & Streamlit 1.57.0, streamlit-drawable-canvas 0.9.3, streamlit-image-coordinates 0.4.1 \\
Suivi d'expériences et tests & MLflow 3.16.0, pytest 9.1.1 \\
\bottomrule
\end{longtable}

\textbf{Expériences réalisées.} Toutes les expériences d'entraînement ont été tracées dans MLflow (hyperparamètres, métriques par époque, checkpoint et journal d'entraînement sauvegardés comme artefacts). Quatre entraînements complets ont été menés avec la même configuration (tableau~\ref{tab:runs}) : le run de référence, puis trois graines différentes pour l'étude de robustesse et l'ensemble de modèles (Chapitre 4).

\begin{longtable}{L{4.6cm} R{1.8cm} R{1.6cm} R{2.2cm} R{3.2cm}}
\caption{Entraînements tracés dans MLflow (encodeur ResNet34, images 384 $\times$ 384, 68 images d'entraînement / 14 de validation).}
\label{tab:runs} \\
\toprule
\textbf{Run} & \textbf{img\_size} & \textbf{Époques} & \textbf{Train / Val} & \textbf{Meilleur Dice (val.)} \\
\midrule
\endfirsthead
\toprule
\textbf{Run} & \textbf{img\_size} & \textbf{Époques} & \textbf{Train / Val} & \textbf{Meilleur Dice (val.)} \\
\midrule
\endhead
\texttt{unet\_resnet34\_holdout} & 384 & 60 & 68 / 14 & 0,923 \\
\texttt{unet\_resnet34\_seed42} & 384 & 60 & 68 / 14 & 0,921 \\
\texttt{unet\_resnet34\_seed123} & 384 & 60 & 68 / 14 & 0,923 \\
\texttt{unet\_resnet34\_seed2026} & 384 & 58 (arrêt anticipé) & 68 / 14 & 0,914 \\
\bottomrule
\end{longtable}

\section{Conception et implémentation de la segmentation}

\textbf{Modèle : U-Net avec encodeur ResNet34.} Encodeur pré-entraîné sur ImageNet (bibliothèque \texttt{segmentation\_models\_pytorch}), une seule classe de sortie (plaie / non-plaie), 3 canaux d'entrée (RGB) :

\begin{lstlisting}[language=Python]
def build_model(encoder_name="resnet34", encoder_weights="imagenet"):
    return smp.Unet(
        encoder_name=encoder_name,
        encoder_weights=encoder_weights,
        in_channels=3,
        classes=1,
    )
\end{lstlisting}

\begin{figure}[h]
\centering
\includegraphics[width=0.9\textwidth]{images/figure_unet_resnet34.png}
\caption{Architecture U-Net avec encodeur ResNet34.}
\label{fig:unet_resnet34}
\end{figure}

\textbf{Prétraitement des images.} Un prétraitement optionnel est proposé côté application (débruitage, correction de champ plat / \textit{flat-field}, égalisation d'histogramme adaptative CLAHE), activable par l'utilisateur avant prédiction. À l'entraînement, l'image est redimensionnée via le pipeline d'augmentation de données (voir 3.3), et non par un simple redimensionnement déterministe ; l'évaluation, elle, est effectuée en résolution native : le masque prédit est ré-suréchantillonné à la résolution originale avant comparaison, pour refléter exactement ce que voit l'utilisateur final dans l'application plutôt qu'un score optimiste obtenu à basse résolution.

\textbf{Reconstruction des masques.} Les masques de vérité terrain utilisés pour l'entraînement sont reconstruits à partir des contours tracés à la main sur les photographies (voir 3.3 pour le détail du traitement).

\textbf{Fonction de perte (Dice + BCE).} Combine une perte de recouvrement (Dice, adaptée aux masques très déséquilibrés fond/plaie) et une perte de classification pixel par pixel (entropie croisée binaire) :

\[
\mathcal{L} = \mathcal{L}_{BCE} + (1 - \mathrm{Dice})
\]

\begin{lstlisting}[language=Python]
class DiceBCELoss(torch.nn.Module):
    def __init__(self):
        super().__init__()
        self.bce = torch.nn.BCEWithLogitsLoss()

    def forward(self, logits, target):
        bce = self.bce(logits, target)
        probs = torch.sigmoid(logits)
        pred = probs.flatten(1)
        tgt = target.flatten(1)
        intersection = (pred * tgt).sum(1)
        dice = (2 * intersection + EPS) / (pred.sum(1) + tgt.sum(1) + EPS)
        return bce + (1 - dice.mean())
\end{lstlisting}

\textbf{Entraînement en deux phases} (\texttt{src/train.py}). L'encodeur pré-entraîné est d'abord gelé pendant 5 époques (seul le décodeur apprend, pour ne pas détruire les poids ImageNet avec un signal de gradient encore bruité), puis dégelé avec un taux d'apprentissage réduit d'un facteur 10 pour affiner l'ensemble du réseau.

\textbf{Hyperparamètres.}

\begin{longtable}{p{7.5cm} r}
\caption{Hyperparamètres d'entraînement du modèle U-Net.}
\label{tab:hyperparametres} \\
\toprule
\textbf{Hyperparamètre} & \textbf{Valeur} \\
\midrule
\endfirsthead
\toprule
\textbf{Hyperparamètre} & \textbf{Valeur} \\
\midrule
\endhead
Époques (maximum) & 60 \\
Taille de batch & 4 \\
Taille d'image & 384 $\times$ 384 \\
Taux d'apprentissage initial & 1$\times$10\textsuperscript{-4} \\
Époques encodeur gelé (warmup) & 5 \\
Patience (early stopping sur Dice val.) & 15 \\
Fonction de perte & Dice + BCE (voir ci-dessus) \\
Ordonnanceur de LR & \texttt{ReduceLROnPlateau} \\
\bottomrule
\end{longtable}

\textbf{Early stopping.} L'entraînement s'arrête automatiquement si le Dice sur le jeu de validation ne s'améliore plus pendant 15 époques consécutives, afin d'éviter le surapprentissage sur un jeu d'entraînement de taille modeste (68 images) plutôt que de fixer arbitrairement un nombre d'époques.

Les courbes d'apprentissage, qui permettent de vérifier que le modèle converge sans signe évident de surapprentissage, sont présentées et analysées au Chapitre 4 (section 4.2).

\textbf{Ensemble de modèles et Test-Time Augmentation.} Au-delà du modèle canonique, \texttt{src/ensemble.py} recombine par vote majoritaire pixel par pixel les prédictions de plusieurs modèles indépendants déjà entraînés (voir 1.8.1, 1.9.4), et produit une carte du nombre de modèles en désaccord par pixel, utilisée comme indicateur d'incertitude. La Test-Time Augmentation moyenne les probabilités prédites sur plusieurs versions retournées de l'image avant seuillage. Les deux approches sont évaluées et comparées au Chapitre 4 ; ce chapitre-ci ne décrit que leur implémentation, pas leurs performances.

\section{Préparation et traitement des données}

\begin{itemize}
\item \textbf{Dataset} : 97 images de microscopie de scratch assay (voir 2.4 pour la source et le contexte).
\item \textbf{Génération des masques à partir des contours} : un script dédié reconstruit les masques binaires à partir des contours tracés à la main directement sur les photographies : détection du trait par couleur/teinte, pontage des discontinuités du trait (effacé par endroits à cause de la compression JPEG), fermeture morphologique sur les bords d'image, puis remplissage par \textit{flood-fill} de la zone délimitée.
\item \textbf{Nettoyage} : une image est exclue du jeu de données final, son tracé manuel étant incomplet même après ce traitement (voir 2.4).
\item \textbf{Split train/validation/test} : séparation reproductible (graine fixée) en 68 images d'entraînement, 14 de validation, 14 de test, le jeu de test étant tenu à l'écart de la sélection du meilleur modèle (voir 2.4).
\item \textbf{Augmentation des données} (bibliothèque \texttt{albumentations}, à l'entraînement uniquement) : recadrage aléatoire, retournements, ajout de bruit --- pour compenser la taille réduite du jeu d'entraînement.
\end{itemize}

\begin{figure}[h]
\centering
\includegraphics[width=0.95\textwidth]{images/figure_reconstruction_masque.png}
\caption{Étapes de reconstruction d'un masque binaire à partir du contour tracé à la main.}
\label{fig:reconstruction_masque}
\end{figure}

\section{Composante RAG/LLM}

\textbf{Corpus documentaire.} La base documentaire du RAG regroupe 48 articles scientifiques relatifs à la cicatrisation et au scratch assay, sélectionnés manuellement (voir 1.8.3, 2.6). Les documents sont indexés dans \textbf{ChromaDB}, interrogée par similarité sémantique (\textit{embeddings}) et complétée par une recherche lexicale \textbf{BM25} (\texttt{rank\_bm25}), pour combiner rapprochement par sens et correspondance de termes exacts lors du \textit{retrieval}. Le modèle de langage est servi en local via \textbf{Ollama} (modèle \textbf{Mistral 7B}~\cite{ref40}), sans appel à une API propriétaire (voir 2.3).

Au-delà de ce mécanisme de récupération documentaire, ce point mérite une sous-section dédiée : la fiabilité de la génération par LLM est la contrainte de conception la plus spécifique de ce projet, largement au-delà d'un simple appel à une API de LLM.

\textbf{Architecture à deux appels séparés.} Le résumé scientifique de cinétique est construit par \textbf{deux appels indépendants} au LLM local :
\begin{enumerate}
\item Un premier appel reçoit uniquement les mesures déjà calculées par le pipeline (surface, \%, évolution entre instants --- voir 3.5) --- jamais d'extrait d'article.
\item Un second appel reçoit uniquement des extraits d'articles scientifiques pertinents --- jamais aucune mesure.
\end{enumerate}

\begin{figure}[h]
\centering
\includegraphics[width=0.95\textwidth]{images/figure_rag_architecture.png}
\caption{Architecture RAG/LLM à deux appels séparés.}
\label{fig:rag_architecture}
\end{figure}

Cette séparation empêche directement le LLM de recevoir, dans un même appel, à la fois les mesures expérimentales et les extraits documentaires : la séparation des deux types de contexte est ainsi assurée au niveau de l'architecture, plutôt que simplement recommandée par une consigne de prompt.

\textbf{Calcul en Python, jamais par le LLM.} Toutes les valeurs numériques présentées (moyennes, évolutions entre instants, comparaison entre intervalles) sont calculées par du code Python pur et injectées déjà calculées dans le prompt ; le LLM est chargé de les restituer et de les mettre en forme, sans être responsable de leur calcul. Ce choix résulte d'observations expérimentales répétées au cours du projet : lorsque le calcul ou la comparaison de chiffres était laissé au modèle de langage local utilisé (Mistral 7B), des erreurs apparaissaient régulièrement --- y compris lorsqu'un patron de phrase explicite lui était fourni pour le faire (bornes d'intervalles inversées, signe négatif fabriqué ou omis). La solution retenue n'a donc pas été de renforcer la consigne de prompt outre mesure, mais de déplacer ces opérations hors du modèle : chaque phrase comparative est entièrement pré-composée par du code Python, que le LLM est chargé de reproduire sans la recalculer.

\textbf{Détection de citations fabriquées.} Un contrôle \textit{a posteriori}, purement programmatique, compare chaque titre d'article cité entre guillemets par le LLM aux titres réellement retournés par la recherche documentaire pour cette génération (comparaison approchée par distance d'édition, pour tolérer des variantes mineures de recopie). Toute citation ne correspondant à aucun titre réel est signalée à l'utilisateur.

\textbf{Autres contrôles après génération.} Sur le même principe (lecture du texte déjà généré, sans modifier ni le prompt ni le contexte), d'autres contrôles déterministes sont appliqués avant affichage : remplacement des références par position (« le premier article ») par le vrai titre de l'article ; suppression d'une phrase de consigne interne que le modèle recopie parfois dans le résumé ; correction d'une faute d'accord récurrente dans la phrase de clôture ; et avertissement lorsque le résumé affirme une augmentation de la surface de plaie alors qu'aucune réouverture n'apparaît dans les mesures.

\textbf{Méthodologie de test dédiée.} Chaque modification du prompt ou du contenu de la base documentaire est validée par un test d'isolation : plusieurs générations réelles sont produites avant et après la modification et comparées, plutôt que de supposer qu'une modification du texte d'instruction suffit à changer le comportement du modèle. Cette méthodologie a permis d'identifier et de corriger plusieurs défauts qui n'auraient pas été détectés par une simple relecture du prompt (voir Chapitre 4.5).

Cette méthodologie a notamment révélé, à plusieurs reprises, une limite structurelle de l'approche par prompt : une reformulation jugée meilleure en moyenne lors d'un test d'isolation peut laisser subsister un défaut déjà connu par ailleurs, ou en révéler un nouveau non détecté lors d'une validation précédente. Dans chaque cas, la correction retenue n'a pas été de renforcer indéfiniment la consigne --- au risque de déstabiliser d'autres parties du prompt, comme observé à plusieurs reprises sur ce projet --- mais un contrôle déterministe en aval, cohérent avec le principe déjà énoncé plus haut. Les défauts concrets identifiés par cette méthodologie, ainsi que leur correction, sont présentés au Chapitre 4.5.

\section{Quantification et analyse temporelle}

\textbf{Calcul de la surface.} La surface de la plaie est obtenue en comptant les pixels du masque prédit appartenant à la classe « plaie » (\texttt{wound\_area\_px}, \texttt{src/measure.py}) --- une simple somme sur le masque binaire, sans hypothèse supplémentaire sur la forme de la région.

\textbf{Calibration pixels/cm.} Lorsqu'un ratio pixels/cm est disponible (mesuré manuellement sur un objet de référence visible dans l'image, ou fourni explicitement), la surface en pixels\textsuperscript{2} est convertie en cm\textsuperscript{2} par \texttt{area\_px\_to\_cm2} ; sans cette calibration, seule la surface en pixels\textsuperscript{2} (et le pourcentage de fermeture, sans unité) est disponible.

\textbf{Pourcentage de fermeture.} Pour une série d'images d'un même échantillon prises à différents instants, la fermeture à l'instant $t$ est calculée par rapport à la surface de référence $S_0$ (le premier instant de la série, généralement 0~h) :
$$\text{Fermeture}(t)\ (\%) = 100 \times \frac{S_0 - S_t}{S_0}$$
(\texttt{src/kinetics.py}) --- cette formule est la même que celle documentée dans la littérature (voir 1.2.3), appliquée ici à des surfaces issues de la segmentation automatique plutôt que d'un tracé manuel.

\textbf{Évolution entre les différents instants.} Les images d'une même série sont triées par temps croissant, et la fermeture est calculée pour chaque instant par rapport à cette même référence $S_0$ --- ce qui donne une évolution de la fermeture au fil des instants observés, plutôt qu'une vitesse de fermeture entre deux instants consécutifs (voir 2.1, qui précise cette même distinction).

\textbf{Génération du résumé à partir des valeurs calculées.} Les grandeurs calculées ici (surface, pourcentage de fermeture, évolution temporelle) sont exactement celles transmises au premier appel LLM décrit en 3.4 ; leur génération en résumé scientifique, et les mécanismes de contrôle associés, sont détaillés dans cette même section 3.4 plutôt que répétés ici.

\section{Interface et aspects MLOps}

\textbf{Interface Streamlit.} L'application (\texttt{app.py}) propose deux onglets : « Analyse d'une image » pour la segmentation, la mesure et la correction manuelle d'une image unique, et « Suivi de cicatrisation » pour le calcul de la cinétique de fermeture sur une série d'images et la génération du résumé scientifique assisté par LLM. Un chat interactif, qui mélangeait mesures et extraits d'articles dans un même appel au LLM, figurait dans une version antérieure ; il a été retiré de l'application pour cette raison.

\textbf{Synthèse vocale du résumé.} Une fois le résumé scientifique généré, un bouton « Écouter le résumé » permet de l'entendre lu à voix haute, et un bouton « Arrêter » d'interrompre la lecture. La solution retenue est la synthèse vocale intégrée au navigateur (\textit{Web Speech API}) plutôt qu'une bibliothèque Python : elle ne demande aucune dépendance supplémentaire ni connexion Internet, s'exécute sur le poste de l'utilisateur avec les voix du système (voix française sous Windows), et évite les problèmes de blocage connus des moteurs de synthèse appelés depuis un serveur Streamlit. Avant la lecture, la mise en forme Markdown du résumé (titres, gras, puces) est retirée pour ne pas être prononcée, et le texte est transmis au navigateur sous forme de chaîne encodée, pour qu'un contenu inattendu ne puisse pas être interprété comme du code. Cette fonctionnalité est purement présentationnelle : elle lit le texte déjà affiché, après les contrôles décrits en 3.4, sans modifier ni les mesures ni le résumé. Elle a été vérifiée dans le navigateur : lecture en français, balises Markdown absentes du texte prononcé, et contenu piégé lu littéralement sans erreur. La figure~\ref{fig:capture_resume} montre le bouton sous un résumé réellement généré.

\begin{figure}[htbp]
\centering
\includegraphics[width=0.85\textwidth]{images/capture_resume_ecoute.png}
\caption{Fin d'un résumé scientifique généré dans l'onglet « Suivi de cicatrisation », avec les boutons « Écouter le résumé » et « Arrêter ».}
\label{fig:capture_resume}
\end{figure}

\textbf{Correction manuelle du masque.} Le contour détecté par le modèle n'est pas toujours parfait : l'utilisateur peut dessiner, au pinceau (largeur réglable), une zone à ajouter à la plaie (tracé vert) ou à retirer (tracé magenta) par-dessus le masque prédit --- l'intérieur du tracé est rempli automatiquement pour éviter un contour dédoublé. Une option permet également d'ignorer entièrement la prédiction du modèle et de tracer la plaie à la main (\texttt{streamlit-drawable-canvas}).

Au-delà du script d'entraînement et de l'application Streamlit, le projet met en place une première couche d'outillage MLOps, avec une portée délibérément limitée (voir 2.3, objectifs non fonctionnels, et Chapitre 5, limites) : il ne s'agit pas de construire un pipeline de déploiement continu complet, mais de rendre l'environnement reproductible et les modèles versionnés.

\textbf{Conteneurisation (Docker).} L'application (Streamlit + dépendances Python) est packagée dans une image Docker CPU-only (les wheels PyTorch CPU sont installées explicitement, pour éviter de télécharger les dépendances CUDA inutiles en l'absence de GPU dans l'image). Un service Ollama séparé est orchestré via \texttt{docker-compose.yml}, avec les données volumineuses (images, sorties, index documentaire) montées en volumes plutôt qu'intégrées à l'image. Un \texttt{HEALTHCHECK} interroge l'endpoint de santé natif de Streamlit pour signaler la disponibilité réelle du service, pas seulement le démarrage du conteneur.

\textbf{Registre de modèles (MLflow Model Registry).} Le script d'entraînement peut enregistrer le meilleur checkpoint comme nouvelle version versionnée d'un modèle nommé dans le registre MLflow (option \texttt{--register\_model}, désactivée par défaut, et sans échec de l'entraînement en cas de problème de registre). Une règle de promotion, volontairement simple, compare le Dice de validation de la nouvelle version à celui de la version actuellement marquée comme version de production (via un alias MLflow, pas les « stages » dépréciés) : la nouvelle version n'est promue que si elle obtient un Dice strictement supérieur, ou si aucune version n'est encore en production. Un échec du registre (indisponibilité, permissions) est capturé et journalisé sans jamais interrompre l'entraînement lui-même --- l'enregistrement du modèle reste une fonctionnalité annexe, pas une dépendance critique du pipeline d'apprentissage.

\textbf{Intégration continue (CI).} Le dépôt GitHub exécute, à chaque modification de code, la suite de tests automatisés puis la construction de l'image Docker, suivie d'un test de fumée qui vérifie que les dépendances de l'image s'importent réellement (pas seulement que la construction de l'image réussit).

\begin{figure}[h]
\centering
\includegraphics[width=0.9\textwidth]{images/figure_mlops.png}
\caption{Architecture simplifiée de l'outillage MLOps du projet.}
\label{fig:mlops}
\end{figure}

\textbf{Limites du déploiement.} Il s'agit d'intégration continue au sens strict --- garantir que le code et l'image restent valides --- et non de déploiement continu : aucune publication d'image ni déploiement automatisé n'a lieu à ce stade (voir Chapitre 5, limites).

\section{Conclusion}

Ce chapitre a décrit comment le système présenté au Chapitre 2 a été concrètement construit : l'architecture logicielle et ses modules, la conception et l'implémentation du modèle de segmentation, la préparation des données d'entraînement, la composante RAG/LLM et ses mécanismes de fiabilité, le calcul des grandeurs de quantification et de suivi temporel, ainsi que l'interface utilisateur et les premiers éléments d'outillage MLOps.

Conformément à la distinction déjà posée aux chapitres précédents, les résultats chiffrés obtenus avec cette implémentation --- comparaison du modèle de segmentation à la baseline classique, apport de l'ensemble de modèles et de la TTA, performance du RAG --- ne sont pas présentés ici : ils font l'objet du chapitre suivant, consacré à l'évaluation et à l'expérimentation.

\chapter{Évaluation et expérimentation}

\section{Métriques d'évaluation utilisées}

Pour une tâche de segmentation binaire, l'accuracy pixel par pixel est peu informative (les masques sont très déséquilibrés : la plaie occupe une faible part de l'image) et n'a volontairement pas été retenue comme métrique principale. Les métriques utilisées sont :

\begin{itemize}
\item \textbf{Dice} (équivalent du F1-score pour des masques binaires) et \textbf{IoU} (Intersection-over-Union) --- métriques de recouvrement, les plus pertinentes pour ce projet.
\item \textbf{Precision} et \textbf{Recall} --- pour caractériser le type d'erreur dominant (sur-segmentation vs sous-segmentation).
\item \textbf{Distance de Hausdorff normalisée} (rapportée à la diagonale de l'image, pour rester comparable entre images de résolutions différentes) --- métrique de pire cas, sensible à un seul point de contour mal placé.
\item \textbf{Erreur de surface} (absolue en px\textsuperscript{2} et relative en \%) --- la métrique la plus directement interprétable pour un biologiste, puisqu'elle porte sur la grandeur finalement utilisée (surface de plaie).
\end{itemize}

\section{Résultats obtenus}
\textbf{Qualité de l'entraînement : le modèle apprend-il correctement ?} Les figures~\ref{fig:loss_curves} et~\ref{fig:dice_curves} montrent l'évolution de la perte et du Dice sur les jeux d'entraînement et de validation (le jeu de test n'est pas suivi époque par époque : il n'est évalué qu'une fois, sur le modèle retenu). La convergence est régulière, et les courbes ne montrent pas de signe évident de surapprentissage sur la période observée : le Dice de validation ne décroche pas de celui d'entraînement. Il lui est même supérieur tout au long de l'entraînement, ce qui est attendu ici : l'augmentation de données forte appliquée à l'entraînement (recadrage aléatoire, retournements, bruit) rend la tâche plus difficile que sur la validation, qui ne subit qu'un redimensionnement. Le meilleur checkpoint (époque 54, Dice de validation 0,923) est retenu selon ce critère.

\begin{figure}[htbp]
\centering
\includegraphics[width=0.8\textwidth]{images/figure3_loss_curves.png}
\caption{Évolution de la perte (Dice + BCE) sur les jeux d'entraînement et de validation.}
\label{fig:loss_curves}
\end{figure}

\begin{figure}[htbp]
\centering
\includegraphics[width=0.8\textwidth]{images/figure4_dice_curves.png}
\caption{Évolution du Dice sur les jeux d'entraînement et de validation.}
\label{fig:dice_curves}
\end{figure}

\textbf{Comparaison U-Net vs baseline classique} (14 images de test, jamais vues à l'entraînement) :

\begin{footnotesize}
\setlength{\tabcolsep}{3pt}
\begin{longtable}{L{3.0cm} R{1.1cm} R{1.1cm} R{1.4cm} R{1.1cm} R{1.7cm} R{1.9cm} R{1.8cm}}
\caption{Comparaison des métriques entre le U-Net et la baseline classique sur le jeu de test.}
\label{tab:unet_vs_baseline} \\
\toprule
\textbf{Méthode} & \textbf{Dice} & \textbf{IoU} & \textbf{Precision} & \textbf{Recall} & \textbf{Hausdorff norm.} & \textbf{MAE surface (px\textsuperscript{2})} & \textbf{Erreur surface rel.} \\
\midrule
\endfirsthead
\toprule
\textbf{Méthode} & \textbf{Dice} & \textbf{IoU} & \textbf{Precision} & \textbf{Recall} & \textbf{Hausdorff norm.} & \textbf{MAE surface (px\textsuperscript{2})} & \textbf{Erreur surface rel.} \\
\midrule
\endhead
Texture + Otsu (baseline) & 0,576 & 0,448 & 0,558 & 0,817 & 0,204 & 210 243 & 404,9 \% \\
\textbf{U-Net + ResNet34} & \textbf{0,882} & \textbf{0,796} & \textbf{0,852} & 0,929 & 0,323 & \textbf{8 538} & \textbf{15,9 \%} \\
\bottomrule
\end{longtable}
\end{footnotesize}

\begin{figure}[htbp]
\centering
\includegraphics[width=\textwidth]{images/figure6_comparaison_methodes.png}
\caption{Comparaison des quatre méthodes sur le jeu de test : métriques de recouvrement (à gauche, plus haut = meilleur) et erreur de surface relative (à droite, plus bas = meilleur). Les résultats de l'ensemble sont exploratoires (voir la réserve méthodologique ci-dessous).}
\label{fig:comparaison_methodes}
\end{figure}

La figure~\ref{fig:comparaison_methodes} résume visuellement cette comparaison, TTA et ensemble compris. Le U-Net domine largement sur les métriques de recouvrement et surtout sur l'erreur de surface (15,9 \% contre 405 \% pour la baseline). Le Hausdorff normalisé du U-Net est en revanche plus élevé (moins bon) que celui de la baseline malgré un Dice très supérieur --- un résultat contre-intuitif rapporté sans le masquer, analysé ci-dessous (Test-Time Augmentation et robustesse inter-seed). Ce résultat montre qu'une métrique prise isolément ne suffit pas à caractériser la qualité d'une segmentation : le Dice et l'IoU décrivent principalement le recouvrement global des régions, tandis que la distance de Hausdorff est particulièrement sensible à des écarts localisés du contour.

\textbf{Test-Time Augmentation (TTA)} --- moyenner les probabilités prédites sur 4 versions retournées de l'image avant seuillage fait chuter le Hausdorff normalisé de 0,323 à 0,043, suggérant que le Hausdorff élevé provient en partie d'artefacts de contour isolés plutôt que d'une dégradation globale du recouvrement. Ce gain n'est cependant pas généralisable à toutes les métriques : le Dice reste quasi inchangé (0,882 $\to$ 0,886) et l'erreur de surface se dégrade légèrement (15,9 \% $\to$ 16,9 \%). La TTA améliore donc spécifiquement la stabilité du contour, sans réduction de l'erreur de surface --- c'est l'ensemble de modèles (ci-dessous) qui obtient la meilleure performance globale parmi les variantes évaluées.

\textbf{Ensemble de 4 modèles + carte de confiance} --- les 4 modèles indépendants déjà entraînés (le run canonique + les 3 graines de l'étude de robustesse) sont combinés par vote majoritaire pixel par pixel ; le nombre de modèles en désaccord à chaque pixel fournit une carte de confiance exploitable directement par l'utilisateur.

\textbf{Limite méthodologique importante.} Les graines 123 et 2026 utilisées ici reposent, comme détaillé dans « Robustesse inter-seed » ci-dessous, sur un split train/validation/test régénéré indépendamment du split canonique. En pratique, chacune des 14 images du jeu de test canonique a été vue, en entraînement ou en validation, par le modèle seed 123, par le modèle seed 2026, ou par les deux --- aucune des 14 images n'est restée non vue par ces deux modèles simultanément. Le tableau ci-dessous ne peut donc pas être interprété comme une performance de test indépendante au même titre que la comparaison U-Net/baseline ci-dessus : il documente le comportement observé de cette combinaison de modèles sur cet échantillon précis, à titre exploratoire, et non un gain de performance validé sur données non vues.

\begin{longtable}{L{5cm} R{3cm} R{5cm}}
\caption{Comparaison entre le modèle unique et l'ensemble de 4 modèles (résultat exploratoire, voir réserve méthodologique ci-dessus).}
\label{tab:ensemble} \\
\toprule
\textbf{Métrique} & \textbf{Modèle unique} & \textbf{Ensemble (4 modèles, résultat exploratoire)} \\
\midrule
\endfirsthead
\toprule
\textbf{Métrique} & \textbf{Modèle unique} & \textbf{Ensemble (4 modèles, résultat exploratoire)} \\
\midrule
\endhead
Dice & 0,882 & 0,893 \\
IoU & 0,796 & 0,812 \\
Precision & 0,852 & 0,881 \\
Recall & 0,929 & 0,915 \\
Hausdorff normalisé & 0,323 & 0,038 \\
Erreur de surface relative & 15,9 \% & 12,2 \% \\
\bottomrule
\end{longtable}

Sous cette réserve, la carte de confiance reste une fonctionnalité à part entière de l'application, indépendamment de la validité du chiffre de Dice ci-dessus : elle signale à l'utilisateur les zones où les modèles de l'ensemble ne s'accordent pas, quelle que soit la performance exacte mesurée sur cet échantillon. La corrélation observée entre le pourcentage de zone en désaccord entre modèles et le Dice par image (-0,751) doit être lue avec la même réserve méthodologique, puisqu'elle est calculée à partir de ce même jeu de prédictions. La figure~\ref{fig:incertitude_dice} représente cette relation image par image : les images au-delà du seuil d'alerte de l'application (25~\%) sont bien celles dont le Dice est le plus faible.
\begin{figure}[htbp]
\centering
\includegraphics[width=0.8\textwidth]{images/figure9_incertitude_vs_dice.png}
\caption{Zone incertaine (désaccord entre les 4 modèles) en fonction du Dice de l'ensemble, sur les 14 images de test (r = -0,751, résultat exploratoire).}
\label{fig:incertitude_dice}
\end{figure}


\textbf{Illustration concrète (deux cas réels du jeu de test, à titre exploratoire).} Pour rendre ce chiffre tangible, deux images de test aux profils opposés sont comparées ci-dessous (carte de confiance : vert = les 4 modèles d'accord, rouge = désaccord partiel) :

\begin{longtable}{L{4cm} R{4.5cm} R{5.5cm}}
\caption{Comparaison des indicateurs de confiance entre les images 6 et 91.}
\label{tab:confiance_illustration} \\
\toprule
 & \textbf{Image 6 (faible incertitude)} & \textbf{Image 91 (forte incertitude)} \\
\midrule
\endfirsthead
\toprule
 & \textbf{Image 6 (faible incertitude)} & \textbf{Image 91 (forte incertitude)} \\
\midrule
\endhead
Zone incertaine & 7,1 \% & \textbf{46,5 \%} ($>$ seuil d'alerte de l'application, 25 \%) \\
Dice & 0,951 & 0,824 \\
Erreur de surface & 5,0 \% & 21,4 \% \\
\bottomrule
\end{longtable}

\begin{figure}[htbp]
\centering
\begin{minipage}{0.48\textwidth}
\centering
\includegraphics[width=\textwidth]{images/figure5a_confiance_faible_incertitude_id6.png}
\caption{Image 6 : faible désaccord entre modèles.}
\label{fig:confiance_faible}
\end{minipage}
\hfill
\begin{minipage}{0.48\textwidth}
\centering
\includegraphics[width=\textwidth]{images/figure5b_confiance_forte_incertitude_id91.png}
\caption{Image 91 : fort désaccord entre modèles.}
\label{fig:confiance_forte}
\end{minipage}
\end{figure}

Sur l'image 6, une plaie franche et continue : les 4 modèles s'accordent presque partout, la zone incertaine reste sous les 10 \% et le Dice dépasse 0,95. Sur l'image 91, une plaie quasi refermée et fragmentée en plusieurs îlots disjoints : les modèles divergent sur près de la moitié de la zone détectée, l'application déclenche son avertissement de vérification manuelle, et le Dice chute à 0,824 --- exactement la relation statistique mesurée ci-dessus, observée ici sur un exemple concret plutôt que sur une moyenne. Ce contraste illustre l'utilité pratique de l'indicateur pour le biologiste : il ne se contente pas d'annoncer une surface, il signale aussi quand cette surface mérite une vérification visuelle avant d'être réutilisée.

\begin{figure}[htbp]
\centering
\includegraphics[width=0.85\textwidth]{images/capture_confiance_zoom.png}
\caption{Onglet « Analyse d'une image » de l'application, estimation de confiance activée, sur l'image 91 : carte d'accord entre les 4 modèles (vert : accord total, rouge : désaccord partiel), zone incertaine affichée à 47 \% (46,5 \% arrondi) et avertissement invitant à vérifier le contour.}
\label{fig:capture_confiance}
\end{figure}

La figure~\ref{fig:capture_confiance} montre ce même indicateur tel que le biologiste le voit dans l'application.

\textbf{Robustesse inter-seed} --- le même protocole d'entraînement a été répété sur 3 graines aléatoires indépendantes pour vérifier que le résultat ne tient pas à un split ou une initialisation favorable. La graine ne contrôle pas uniquement l'initialisation du modèle : elle contrôle aussi, via \texttt{src/split\_data.py}, la génération du split train/validation/test lui-même. Les graines 123 et 2026 régénèrent donc un split entièrement différent du split canonique (seule la graine 42 coïncide avec le split canonique, ce dernier ayant été généré avec cette même valeur par défaut) ; les quatre modèles comparés ici (canonique compris) sont par ailleurs quatre entraînements indépendants, jamais une réutilisation du même checkpoint :

\begin{small}
\setlength{\tabcolsep}{4pt}
\begin{longtable}{L{3.2cm} R{2.6cm} R{2.6cm} R{2.8cm} R{2.8cm}}
\caption{Résultats de l'étude de robustesse inter-seed.}
\label{tab:robustesse_seed} \\
\toprule
\textbf{Seed} & \textbf{Dice} & \textbf{IoU} & \textbf{Erreur de surface} & \textbf{Hausdorff normalisé} \\
\midrule
\endfirsthead
\toprule
\textbf{Seed} & \textbf{Dice} & \textbf{IoU} & \textbf{Erreur de surface} & \textbf{Hausdorff normalisé} \\
\midrule
\endhead
42 & 0,880 & 0,791 & 16,1 \% & 0,083 \\
123 & 0,903 & 0,828 & 7,2 \% & 0,131 \\
2026 & 0,892 & 0,809 & 8,8 \% & 0,032 \\
\textbf{moyenne $\pm$ écart-type} & \textbf{0,892 $\pm$ 0,009} & \textbf{0,809 $\pm$ 0,015} & \textbf{10,7 \% $\pm$ 3,9 \%} & \textbf{0,082 $\pm$ 0,040} \\
\bottomrule
\end{longtable}
\end{small}

\begin{figure}[htbp]
\centering
\includegraphics[width=\textwidth]{images/figure8_robustesse_seeds.png}
\caption{Dice moyen (à gauche) et distance de Hausdorff normalisée moyenne (à droite) du run de référence et des trois graines, sur le jeu de test de chaque run.}
\label{fig:robustesse_seeds}
\end{figure}

La figure~\ref{fig:robustesse_seeds} met en évidence le contraste entre les deux métriques. Le Dice reste relativement stable malgré les variations de graine et de découpage des données observées dans cette expérimentation. Ce résultat suggère une certaine robustesse de la performance aux variations du protocole de séparation et de l'initialisation. Il ne permet toutefois pas de conclure à une robustesse générale, notamment en l'absence de validation externe. Le Hausdorff, en revanche, varie fortement (presque $\times 4$ entre graines), confirmant qu'il ne doit pas être lu comme un score de qualité stable pour cette tâche.

\section{Comparaison de différents modèles/paramètres et significativité statistique}

Au-delà de la comparaison de moyennes, un test des rangs signés de Wilcoxon (non paramétrique, adapté à un petit échantillon) a été appliqué aux paires \textbf{(U-Net, baseline classique)} sur le Dice et sur l'IoU de chacune des 14 images de test (comparaison appariée image par image, script \texttt{src/statistical\_comparison.py}), complété par un intervalle de confiance bootstrap (10 000 ré-échantillonnages) sur la différence moyenne :

\begin{footnotesize}
\setlength{\tabcolsep}{3pt}
\begin{longtable}{L{1.4cm} R{1.5cm} R{1.7cm} R{1.5cm} R{2.0cm} R{1.7cm} R{3.2cm}}
\caption{Test de Wilcoxon et intervalle de confiance bootstrap (U-Net vs baseline).}
\label{tab:wilcoxon} \\
\toprule
\textbf{Métrique} & \textbf{Moy. U-Net} & \textbf{Moy. baseline} & \textbf{Diff. moy.} & \textbf{U-Net meilleur sur} & \textbf{p (Wilcoxon)} & \textbf{IC bootstrap 95\%} \\
\midrule
\endfirsthead
\toprule
\textbf{Métrique} & \textbf{Moy. U-Net} & \textbf{Moy. baseline} & \textbf{Diff. moy.} & \textbf{U-Net meilleur sur} & \textbf{p (Wilcoxon)} & \textbf{IC bootstrap 95\%} \\
\midrule
\endhead
Dice & 0,882 & 0,576 & +0,307 & 14/14 images & 0,000122 & [+0,183 ; +0,448] \\
IoU & 0,796 & 0,448 & +0,348 & 14/14 images & 0,000122 & [+0,236 ; +0,470] \\
\bottomrule
\end{longtable}
\end{footnotesize}

\begin{figure}[htbp]
\centering
\includegraphics[width=\textwidth]{images/figure7_dice_par_image.png}
\caption{Dice du U-Net et de la baseline pour chacune des 14 images de test (triées par Dice du U-Net). Chaque segment relie les deux scores d'une même image.}
\label{fig:dice_par_image}
\end{figure}

La figure~\ref{fig:dice_par_image} montre ce résultat image par image : aucun segment ne descend, le U-Net est au-dessus de la baseline pour chaque image. Sur le Dice et sur l'IoU, le U-Net obtient un score supérieur à la baseline classique sur les 14 images de test sans exception, ce qui produit le p le plus faible que ce test puisse atteindre à cette taille d'échantillon (largement $<$ 0,05) : \textbf{la supériorité du U-Net sur la baseline classique} (et non sur l'ensemble de modèles ou une autre variante) est statistiquement significative sur cet échantillon. Cette comparaison ne porte que sur U-Net vs baseline ; les gains de la TTA (4.2) sont rapportés comme des moyennes sur le même jeu de test, sans test de significativité dédié à ce stade, et ceux de l'ensemble de modèles (4.2) le sont à titre exploratoire compte tenu de la limite méthodologique déjà signalée sur les splits. Réserve méthodologique : un échantillon de 14 images reste petit pour généraliser au-delà de ce jeu de test précis --- réserve que l'étude de robustesse inter-seed (4.2, Dice = 0,892 $\pm$ 0,009 sur 3 graines) vient partiellement contrebalancer, en montrant que le résultat ne tient pas à un split ou une initialisation favorable.

Les comparaisons de paramètres/configurations réalisées sont : modèle unique vs ensemble de 4 modèles (4.2), avec vs sans TTA (4.2), et robustesse sur 3 graines d'initialisation différentes (4.2) --- un seul jeu d'hyperparamètres d'entraînement (encodeur, taille d'image, taux d'apprentissage) a été exploré à ce stade (voir Chapitre 5, perspectives).

\section{Analyse critique des résultats}

\begin{figure}[htbp]
\centering
\includegraphics[width=\textwidth]{images/figure1_exemples.png}
\caption{Quatre cas représentatifs du jeu de test (image originale, vérité terrain, U-Net, baseline) : bonne segmentation (image 39), cas difficile (image 50), plaie très petite (image 93) et échec de la baseline sous éclairage atypique (image 76).}
\label{fig:exemples}
\end{figure}

\textbf{Cas qualitatifs représentatifs} (figure~\ref{fig:exemples}, quatre cas : image originale / vérité terrain / U-Net / baseline) : sur un cas de bonne segmentation, les deux méthodes suivent le contour mais la baseline ajoute du bruit ; sur un cas difficile, le U-Net reste précis quand la baseline se fragmente ; sur un cas d'échec de la baseline (éclairage atypique, fond grisâtre plutôt que rose), le U-Net suit correctement le contour quand la baseline dérive et trace presque tout le pourtour du champ --- illustrant concrètement la fragilité d'une méthode de seuillage classique face à des conditions d'acquisition variables.

\textbf{Analyse d'erreur quantitative par catégorie} --- pour vérifier avec des chiffres (pas seulement sur les cas illustrés) l'hypothèse de fragilité de la baseline à l'éclairage, les images de test ont été croisées avec un critère objectif de luminosité (z-score par rapport au dataset complet) :

\begin{longtable}{l r r}
\caption{Comparaison du Dice par catégorie d'éclairage.}
\label{tab:eclairage} \\
\toprule
\textbf{Éclairage (test, n)} & \textbf{Dice U-Net} & \textbf{Dice baseline} \\
\midrule
\endfirsthead
\toprule
\textbf{Éclairage (test, n)} & \textbf{Dice U-Net} & \textbf{Dice baseline} \\
\midrule
\endhead
Normal (12) & 0,884 & 0,656 \\
Atypique (2) & 0,874 & \textbf{0,092} \\
\bottomrule
\end{longtable}

Sur les images à éclairage atypique, le U-Net reste stable quand la baseline s'effondre. Ces résultats sont compatibles avec l'hypothèse d'une plus grande sensibilité de la baseline aux variations d'éclairage ; toutefois, compte tenu du faible nombre d'images atypiques ($n=2$), cette observation doit être considérée comme une tendance exploratoire plutôt que comme une conclusion statistiquement établie. \textbf{Avertissement méthodologique explicite} : avec seulement 14 images de test, chaque catégorie ne contient que 1 à 13 images ; ce tableau sert à confirmer une tendance déjà observée qualitativement, pas à établir un résultat statistique à part entière.

\textbf{Cas limite documenté} : l'image obtenant le pire score des deux méthodes correspond à la plus petite surface de plaie de tout le dataset (rang 1/97, 1,3 \% du champ). Sur une structure aussi fine, quelques pixels d'écart de contour suffisent à faire chuter fortement des métriques proportionnelles à la surface (Dice/IoU) --- une limite connue et mesurable sur les plaies quasi refermées, partagée par les deux méthodes, et non un échec de généralisation isolé. La figure~\ref{fig:surface_dice} montre cette tendance sur les 14 images de test : le Dice du U-Net diminue pour les plaies de très petite surface relative. Avec seulement 14 points, il s'agit d'un signal cohérent à garder en tête pour les plaies quasi refermées, pas d'une relation statistique démontrée.

\begin{figure}[htbp]
\centering
\includegraphics[width=0.75\textwidth]{images/figure2_surface_vs_dice.png}
\caption{Dice du U-Net en fonction de la surface relative de la plaie (\% du champ), sur les 14 images de test.}
\label{fig:surface_dice}
\end{figure}

\textbf{Mise en regard avec la littérature (Chapitre 1)} : le Dice de 0,882 obtenu ici, sur seulement 97 images, est cohérent avec un comparateur indépendant sur une modalité d'imagerie proche (Dice 0,876, cellules souches mésenchymateuses \cite{ref3}) mais nettement inférieur à un comparateur direct sur scratch assay disposant de plus de 400 images (Dice 0,958--0,968 \cite{ref2}) --- l'écart le plus plausible étant la taille du jeu d'entraînement, pas un défaut d'architecture (voir Chapitre 5). Ces comparaisons doivent néanmoins être interprétées avec prudence : il ne s'agit pas d'une comparaison expérimentale directe, les jeux de données, protocoles d'annotation et conditions d'imagerie différant d'une étude à l'autre.

\textbf{Synthèse : le modèle est-il satisfaisant ?} Le tableau~\ref{tab:verdict} répond à cette question critère par critère, à partir des résultats de ce chapitre.

\begin{small}
\begin{longtable}{L{3.6cm} L{6.6cm} L{3.6cm}}
\caption{Évaluation du modèle U-Net + ResNet34, critère par critère.}
\label{tab:verdict} \\
\toprule
\textbf{Critère} & \textbf{Résultat mesuré} & \textbf{Verdict} \\
\midrule
\endfirsthead
\toprule
\textbf{Critère} & \textbf{Résultat mesuré} & \textbf{Verdict} \\
\midrule
\endhead
Apprentissage & Dice de validation 0,923 ; courbes d'entraînement et de validation parallèles (figures~\ref{fig:loss_curves} et~\ref{fig:dice_curves}) & Pas de signe évident de surapprentissage \\
Gain par rapport à la méthode classique & Dice 0,882 contre 0,576 ; meilleur sur 14 images sur 14 ; $p = 0,000122$ & Satisfaisant : gain net et significatif \\
Mesure de surface (objectif du projet) & Erreur relative 15,9~\% contre 404,9~\% pour la baseline & Suffisant pour suivre une tendance ; à vérifier sur les images signalées \\
Stabilité & Dice 0,892 $\pm$ 0,009 sur 3 graines & Satisfaisant \\
Robustesse à l'éclairage & Dice 0,874 sur éclairage atypique contre 0,884 en conditions normales & Satisfaisant, mais sur 2 images seulement \\
Type d'erreur & Rappel 0,929 supérieur à la précision 0,852 & Légère sur-segmentation \\
Plaies quasi refermées & Dice 0,651 sur la plus petite plaie (image 93) & Point faible \\
Précision du contour & Hausdorff normalisé 0,323, très variable selon la graine & Partiel (amélioré par TTA et ensemble) \\
Niveau par rapport à la littérature & Proche d'un comparateur indépendant (0,876), en dessous d'une étude à plus de 400 images (0,958--0,968) & Cohérent pour 97 images \\
Généralisation à un autre laboratoire & Non évaluée & Inconnu \\
\bottomrule
\end{longtable}
\end{small}

\textbf{Conclusion de l'évaluation.} Le modèle n'est pas présenté comme excellent dans l'absolu. Les résultats suggèrent qu'il est suffisamment performant pour automatiser la segmentation et le suivi de la fermeture dans les conditions expérimentales étudiées : il obtient un Dice de 0,882 contre 0,576 pour la méthode classique, réduit fortement l'erreur de surface (environ 16~\% contre plus de 400~\%) et conserve des performances proches d'un entraînement à l'autre (Dice 0,892 $\pm$ 0,009 sur 3 graines). En revanche, sa généralisation à d'autres microscopes et laboratoires reste à valider, et les plaies quasi refermées ou les cas signalés par la carte de confiance demandent une vérification visuelle.

\section{Évaluation du composant RAG/LLM}

\textbf{Fiabilité de la génération par LLM (évaluation qualitative)} --- au-delà des métriques de segmentation, le projet a développé une méthodologie de test spécifique pour la partie LLM/RAG : tests de régression exécutant de vraies générations contre le modèle local et vérifiant que le texte produit contient exactement les valeurs calculées par le pipeline et aucune valeur fabriquée connue. Cette méthodologie a permis de détecter, au cours du développement, plusieurs défauts non triviaux (confusion entre le taux d'un intervalle et la vitesse moyenne globale, signe négatif fabriqué sur une vitesse pourtant positive, bornes d'intervalle inversées) --- chacun corrigé non pas en renforçant la consigne donnée au modèle, mais en supprimant sa liberté de calcul sur le point concerné (voir Chapitre 3.4). C'est une contribution méthodologique du projet à part entière : la fiabilité d'un système LLM+RAG dans un contexte scientifique ne peut pas être supposée à partir de la qualité apparente d'un prompt, elle doit être mesurée par génération réelle répétée.

Deux défauts supplémentaires, identifiés lors de tests d'usage réel plus tardifs, illustrent cette même limite structurelle de l'approche par prompt. D'une part, une reformulation du prompt de mise en contexte bibliographique --- jugée meilleure en moyenne en test d'isolation --- n'a pas éliminé un défaut déjà connu par ailleurs : citer une source par sa position (« la première étude ») plutôt que par son titre. La correction retenue a été, comme pour la détection de citations fabriquées (voir 3.4), un contrôle programmatique après coup, remplaçant systématiquement toute référence positionnelle par le vrai titre de l'article avant affichage à l'utilisateur. D'autre part, un autre test a révélé que le paragraphe de mise en contexte présentait parfois des conclusions biologiques précises (« les cellules migrent de façon coordonnée ») comme des faits généraux acquis, alors que les extraits ne font que proposer des méthodes pour étudier ce type de comportement --- avec un risque concret de laisser croire que l'application mesure elle-même la trajectoire ou la directionnalité des cellules, alors qu'elle ne mesure que la surface de la plaie et son évolution temporelle. Le prompt a été ajusté en conséquence (décrire ce que les travaux \textit{proposent}, jamais ce qu'ils \textit{établissent} comme fait biologique général) et retesté en isolation, avec une amélioration nette sur ce défaut précis et sans régression nouvelle par rapport aux limites déjà documentées de ce modèle local.

\textbf{Évaluation du retrieval documentaire (RAG)} --- un benchmark reproductible de questions représentatives associées manuellement à leur(s) article(s) pertinent(s) donne un Recall@3 de 69,2 \% et un MRR@10 de 0,561 sur le corpus actuel (48 articles, 13 questions évaluées) : la recherche sémantique retrouve le bon article dans le top 3 pour environ sept questions testées sur dix, avec une sensibilité documentée à la formulation exacte de la question plutôt qu'à un simple écart de langue français/anglais (voir Chapitre 5, perspectives).

\section{Conclusion}

Ce chapitre a permis d'évaluer expérimentalement les différentes composantes du système développé. Sur le jeu de test composé de 14 images, le modèle U-Net avec encodeur ResNet34 obtient un Dice de 0,882, contre 0,576 pour la méthode classique utilisée comme baseline --- une différence statistiquement significative (test de Wilcoxon, $p = 0,000122$), mais limitée à cet échantillon. L'erreur relative sur la surface est également fortement réduite, passant de 404,9 \% pour la baseline à 15,9 \% pour le modèle profond.

L'étude de plusieurs configurations a exploré l'apport d'un ensemble de quatre modèles, qui obtient un Dice de 0,893 sur le jeu de test canonique --- un résultat à interpréter à titre exploratoire plutôt que comme une performance de test indépendante, deux des quatre modèles ayant été entraînés sur un split qui chevauche largement ce même jeu de test (voir 4.2). La carte de désaccord entre les modèles reste néanmoins une fonctionnalité utile pour identifier les images nécessitant une vérification supplémentaire, indépendamment de cette réserve. L'étude de robustesse inter-seed suggère par ailleurs une certaine stabilité du Dice face aux variations de split et d'initialisation, sans permettre de conclure à une robustesse générale en l'absence de validation externe ; la distance de Hausdorff, elle, reste nettement plus variable et ne doit pas être lue comme un score de qualité stable pour cette tâche.

L'évaluation du composant RAG/LLM a montré que le système permet de retrouver une partie significative des documents pertinents du corpus scientifique, avec un Recall@3 de 69,2 \% et un MRR@10 de 0,561. Les tests de régression développés spécifiquement pour cette partie du projet ont par ailleurs permis d'identifier et de corriger plusieurs comportements indésirables dans la génération des réponses, sans jamais pouvoir supposer qu'un prompt seul garantit un comportement fiable.

Ces résultats montrent que le prototype répond aux objectifs fixés en matière de segmentation, de quantification et d'assistance à l'interprétation. Ils doivent toutefois être interprétés en tenant compte de la taille limitée du jeu de données, de l'absence de validation externe et des limites propres à certaines métriques et au système RAG. Ces limites, ainsi que les principales perspectives d'amélioration, sont présentées dans le chapitre suivant.

\chapter{Conclusion et perspectives}

\section{Bilan et principales réalisations}

Ce projet visait à automatiser la segmentation et la quantification de la zone de plaie dans des images de scratch assay à l'aide du Deep Learning, tout en évaluant sa fiabilité par rapport à une méthode classique de traitement d'image plutôt que de supposer a priori sa supériorité. Les réalisations du projet peuvent se regrouper en quatre contributions :

\textbf{Vision par ordinateur.} Un pipeline complet de segmentation automatique de scratch assay (U-Net + ResNet34), évalué rigoureusement sur un jeu de test indépendant et jamais vu à l'entraînement, avec une différence statistiquement significative entre le U-Net et la méthode classique de traitement d'image sur le jeu de test évalué (test de Wilcoxon, p = 0,000122). Une étude inter-seed a permis d'examiner la stabilité des performances face aux variations de graine et de découpage des données. Une méthode d'ensemble de modèles a également été développée afin de fournir une carte de confiance par pixel exploitable directement par l'utilisateur final, illustrée concrètement sur des cas réels --- la mesure quantitative de cet ensemble sur le jeu de test canonique restant toutefois exploratoire, deux des quatre modèles ayant été entraînés sur un split chevauchant ce même jeu de test (Chapitre 4.2).

\textbf{Quantification.} Transformation du masque de segmentation en grandeurs directement utiles au biologiste (surface, pourcentage de fermeture, évolution entre les instants observés), calculées uniquement en Python et jamais par le LLM (Chapitre 3.5).

\textbf{RAG/LLM.} Une architecture LLM/RAG conçue spécifiquement pour empêcher la fabrication de chiffres ou de citations scientifiques, validée par une méthodologie de test dédiée (génération réelle répétée, pas une simple relecture de prompt).

\textbf{Applicatif et MLOps.} Une application Streamlit complète, en deux onglets, intégrant mesure, suivi temporel, correction manuelle, interprétation assistée par LLM et écoute du résumé par synthèse vocale, packagée dans une image Docker, avec un suivi d'expériences et un registre de modèles MLflow, et une intégration continue sur GitHub Actions.

\section{Difficultés rencontrées}

\begin{itemize}
\item \textbf{Taille réduite du jeu de données} (97 images au total) : a orienté plusieurs choix de conception (transfer learning, augmentation de données forte, réserve méthodologique sur la portée des résultats statistiques).
\item \textbf{Instabilité de la distance de Hausdorff}, initialement difficile à interpréter (le U-Net semblait moins bon que la baseline sur cette seule métrique) : examinée par une double investigation (Test-Time Augmentation, puis étude multi-seed), qui suggère que cette instabilité est fortement liée à la géométrie fine de la plaie et aux écarts localisés de contour plutôt qu'à une dégradation globale du modèle.
\item \textbf{Fiabilité d'un LLM local de taille modeste} (Mistral 7B) sur des tâches de recopie/comparaison numérique : plusieurs itérations ont été nécessaires pour aboutir à une architecture qui élimine structurellement le risque de fabrication, une consigne de prompt seule s'étant révélée insuffisante à plusieurs reprises.
\item \textbf{Sensibilité du retrieval documentaire à la formulation de la question}, en particulier pour des questions de comparaison quantitative (« mon Dice est-il faible comparé à d'autres études ? »), non résolue de façon entièrement satisfaisante à ce stade.
\end{itemize}

\section{Limites}

Au-delà des difficultés déjà résolues au cours du développement (5.2), le système présente des limites intrinsèques qui doivent être prises en compte dans l'interprétation de ses résultats et dans son usage :

\begin{itemize}
\item \textbf{Jeu de données restreint} : 97 images au total, dont seulement 14 dans le jeu de test. Les résultats obtenus, dont la différence entre le U-Net et la baseline est statistiquement significative sur le Dice et l'IoU dans cet échantillon (Chapitre 4.3), ne garantissent pas une généralisation à d'autres laboratoires, microscopes ou protocoles d'acquisition (voir Chapitre 1.7.5 ; Chapitre 4.4, mise en regard avec la littérature).
\item \textbf{Absence de validation externe} : aucune image issue d'un laboratoire ou d'un microscope différent de celui utilisé pour l'entraînement n'a été évaluée à ce stade.
\item \textbf{Évaluation de l'ensemble de modèles à revoir} : deux des quatre modèles de l'ensemble (seeds 123 et 2026) ont été entraînés sur un split train/validation/test régénéré indépendamment du split canonique, si bien qu'aucune des 14 images du test canonique n'est restée non vue par ces deux modèles simultanément. Le Dice de 0,893 rapporté pour l'ensemble (Chapitre 4.2) doit donc être lu comme un résultat exploratoire et non comme une performance de test indépendante ; une évaluation méthodologiquement valide nécessiterait de réentraîner les modèles de l'ensemble sur un split commun.
\item \textbf{Sensibilité de certaines métriques} : la distance de Hausdorff s'est révélée instable et sensible à la géométrie fine de la plaie (Chapitre 4.2), ce qui limite son interprétation isolée comme score de qualité.
\item \textbf{Portée du corpus documentaire du RAG} : 48 articles scientifiques sélectionnés manuellement, avec un retrieval sensible à la formulation exacte de la question (Recall@3 de 69,2 \%, Chapitre 4.5) --- un corpus plus large ou mieux structuré par catégorie pourrait améliorer ce résultat.
\item \textbf{Portée du modèle de langage local} : Mistral 7B, servi en local pour rester sans dépendance à une API propriétaire, reste sujet à des erreurs de calcul ou de formulation lorsqu'on le laisse libre --- d'où l'architecture de contrôle décrite au Chapitre 3.4, qui limite ce risque sans l'éliminer totalement pour les aspects non couverts par un contrôle programmatique.
\item \textbf{Portée de la synthèse vocale} : la qualité et la disponibilité de la voix dépendent du navigateur et des voix installées sur le poste (une voix française n'est pas garantie sur tous les systèmes), et la lecture n'est pas exportable sous forme de fichier audio.
\item \textbf{Portée du déploiement} : conteneurisation Docker et intégration continue (tests + build + test de fumée) sont en place, mais aucun déploiement automatisé, aucun monitoring de dérive des données en production et aucune chaîne de déploiement continu (CD) n'ont été mis en œuvre (Chapitre 3.6) --- une portée volontairement limitée à ce stade du projet.
\item \textbf{Portée d'usage} : l'application est un outil de quantification pour la recherche in vitro et n'est explicitement pas destinée à un usage diagnostique ou clinique (voir Chapitre 2.3).
\end{itemize}

\section{Perspectives d'évolution}

\textbf{Données.}
\begin{itemize}
\item \textbf{Validation externe} : évaluer le modèle sur des images de scratch assay d'un autre laboratoire/microscope, pour démontrer une généralisation au-delà des conditions d'acquisition du jeu de données actuel.
\item \textbf{Enrichissement du jeu de données} : élargir le nombre d'images annotées, en particulier sur les cas actuellement sous-représentés (plaies quasi refermées, éclairages atypiques).
\end{itemize}

\textbf{Segmentation.}
\begin{itemize}
\item \textbf{Exploration d'hyperparamètres plus large} : un seul jeu d'hyperparamètres d'entraînement a été tracé à ce stade (encodeur, taille d'image, taux d'apprentissage) ; une recherche plus systématique permettrait de savoir si la configuration actuelle est déjà proche d'un optimum.
\item \textbf{Comparaison architecturale} : confronter U-Net/ResNet34 à d'autres architectures de segmentation légères, pour distinguer l'apport du deep learning en général de celui de cette configuration précise.
\item \textbf{Ensemble de modèles méthodologiquement valide} : réentraîner les modèles de l'ensemble sur un split train/validation/test commun (seule la graine d'initialisation variant d'un modèle à l'autre), afin d'obtenir un Dice d'ensemble réellement mesuré sur des données non vues par aucun des modèles combinés (voir 5.3, limites).
\end{itemize}

\textbf{RAG/LLM.}
\begin{itemize}
\item \textbf{Amélioration du retrieval documentaire} : la recherche hybride mots-clés + embeddings, qui déstabilisait le résumé tant que mesures et extraits partageaient le même appel au LLM, est désormais active pour la partie littérature du résumé, depuis la séparation en deux appels (retestée en isolation, sans fabrication numérique). Un modèle d'embeddings multilingue plus grand reste une piste, à valider par le même double protocole de test (benchmark de retrieval ET test d'isolation de génération).
\item \textbf{Base documentaire recentrée par catégorie} : plutôt qu'un élargissement généraliste du corpus, cibler spécifiquement les catégories les plus utiles à l'interprétation d'un résultat de scratch assay (principe et méthodologie, quantification et analyse d'images, migration cellulaire, prolifération --- pour éviter une lecture simpliste « fermeture = migration », cinétique de fermeture, facteurs expérimentaux pouvant expliquer un résultat inhabituel), pour que chaque référence retrouvée corresponde à une vraie question qu'un biologiste poserait sur son résultat, plutôt qu'une recherche générale sur la cicatrisation.
\end{itemize}

\textbf{MLOps et déploiement.}
\begin{itemize}
\item \textbf{Chaîne de déploiement continu (CD)} : l'intégration continue actuelle s'arrête à la validation de l'image Docker (voir 3.6) ; une chaîne complète irait jusqu'à la publication et au déploiement automatisés, dans un cadre encore à définir (hébergement, gestion des accès).
\item \textbf{Monitoring en production} : suivre les performances du modèle déployé et détecter une éventuelle dérive des données (nouvelles conditions d'acquisition, nouveau microscope) plutôt que de supposer que les performances mesurées au Chapitre 4 restent valables indéfiniment.
\item \textbf{Réentraînement contrôlé} : s'appuyer sur le registre de modèles MLflow déjà en place (voir 3.6) pour comparer objectivement une nouvelle version du modèle, entraînée sur des données enrichies, à la version actuellement en production, avant toute promotion.
\end{itemize}

\section{Conclusion générale}

Ce projet a abouti à un prototype de recherche complet combinant vision par ordinateur, quantification, interprétation assistée par LLM/RAG et premiers éléments d'outillage MLOps, pour automatiser une tâche jusque-là réalisée manuellement. Le gain apporté par le Deep Learning par rapport à une méthode classique a été mesuré plutôt que supposé, et les limites de ce prototype --- taille du jeu de données, absence de validation externe, portée du corpus documentaire et du modèle de langage local, portée volontairement limitée du déploiement --- ont été explicitement documentées plutôt que passées sous silence.

L'application reste un outil de quantification pour la recherche in vitro : elle assiste l'interprétation d'un résultat expérimental, mais ne constitue ni un outil diagnostique ni un outil clinique. Les perspectives présentées dans ce chapitre --- validation externe, enrichissement des données, amélioration du RAG et de la segmentation, développement d'une chaîne MLOps complète --- définissent une trajectoire d'évolution cohérente avec cette portée, plutôt qu'un changement de nature du projet.

% ============================================================
% BIBLIOGRAPHIE
% ============================================================
% À COMPLÉTER : les auteurs ne sont pas tous connus (seuls titre/année sont
% tracés dans la base documentaire du projet pour la plupart des entrées) —
% à récupérer avant mise en forme définitive (style IEEE ou APA selon les
% consignes de l'établissement). Sélection couvrant les catégories utilisées
% au Chapitre 1 ; l'intégralité du corpus (48 articles) est disponible dans
% chatbot/articles_cicatrisation.json.
\renewcommand{\bibname}{Bibliographie}
\begin{thebibliography}{99}
\addcontentsline{toc}{chapter}{Bibliographie}

\bibitem{ref1} Ronneberger, O. et al. « U-Net: Convolutional Networks for Biomedical Image Segmentation ». 2015.
\bibitem{ref2} « An automated in vitro wound healing microscopy image analysis approach utilizing U-net-based deep learning methodology ». 2024.
\bibitem{ref3} « Comparative Study of Deep Transfer Learning Models for Semantic Segmentation of Human Mesenchymal Stem Cell Micrographs ». 2025.
\bibitem{ref4} « TernausNet: U-Net with VGG11 Encoder Pre-Trained on ImageNet for Image Segmentation ». 2018.
\bibitem{ref5} « Embracing imperfect datasets: A review of deep learning solutions for medical image segmentation ». 2020.
\bibitem{ref6} « Metrics reloaded: recommendations for image analysis validation ». 2024.
\bibitem{ref7} « Metrics for evaluating 3D medical image segmentation: analysis, selection, and tool ». 2015.
\bibitem{ref8} « clDice --- A Novel Topology-Preserving Loss Function for Tubular Structure Segmentation ». 2021.
\bibitem{ref9} « Studying Cell Migration (Random and Wound Healing) Parameters with Imaging and MATLAB Analysis ». 2023.
\bibitem{ref10} « Migratory Metrics of Wound Healing: A Quantification Approach for in vitro Scratch Assays ». 2018.
\bibitem{ref11} « Study of Wound Healing Dynamics by Single Pseudo-Particle Tracking in Phase Contrast Images Acquired in Time-Lapse ». 2021.
\bibitem{ref12} « Modeling epithelial wound closure dynamics with AI: A comparative study across cell types ». 2025.
\bibitem{ref13} « The In Vitro Wound-Scratch Assay: Applications, Technical Advances, and Limitations in Wound Healing Research ». 2026.
\bibitem{ref14} « Research Techniques Made Simple: Analysis of Collective Cell Migration Using the Wound Healing Assay ». 2017.
\bibitem{ref15} « TScratch: a novel and simple software tool for automated analysis of monolayer wound healing assays ». 2009.
\bibitem{ref16} « Fiji: an open-source platform for biological-image analysis ». 2012.
\bibitem{ref17} « AIM: A Computational Tool for the Automatic Quantification of Scratch Wound Healing Assays ». 2017.
\bibitem{ref18} « A Robust and Standardized Approach to Quantify Wound Closure Using the Scratch Assay ». 2023.
\bibitem{ref19} Buddenkotte, T. et al. « Calibrating Ensembles for Scalable Uncertainty Quantification in Deep Learning-based Medical Segmentation ». \textit{Computers in Biology and Medicine}, 2023.
\bibitem{ref20} « A Review of Uncertainty Estimation and its Application in Medical Imaging ». 2023.
\bibitem{ref21} Otsu, N. « A Threshold Selection Method from Gray-Level Histograms ». \textit{IEEE Transactions on Systems, Man, and Cybernetics}, 1979.
\bibitem{ref22} He, K. et al. « Deep Residual Learning for Image Recognition ». \textit{CVPR}, 2016.
\bibitem{ref23} Deng, J. et al. « ImageNet: A Large-Scale Hierarchical Image Database ». \textit{CVPR}, 2009.
\bibitem{ref24} Lewis, P. et al. « Retrieval-Augmented Generation for Knowledge-Intensive NLP Tasks ». \textit{NeurIPS}, 2020.

\bibitem{ref25} Paszke, A. et al. « PyTorch: An Imperative Style, High-Performance Deep Learning Library ». \textit{Advances in Neural Information Processing Systems (NeurIPS)}, 2019.
\bibitem{ref26} Iakubovskii, P. « Segmentation Models Pytorch ». Dépôt GitHub, 2019. \url{https://github.com/qubvel/segmentation_models.pytorch}
\bibitem{ref27} Buslaev, A. et al. « Albumentations: Fast and Flexible Image Augmentations ». \textit{Information}, vol. 11, n\textsuperscript{o} 2, 2020.
\bibitem{ref28} Bradski, G. « The OpenCV Library ». \textit{Dr. Dobb's Journal of Software Tools}, 2000.
\bibitem{ref29} Harris, C. R. et al. « Array programming with NumPy ». \textit{Nature}, vol. 585, 2020.
\bibitem{ref30} Virtanen, P. et al. « SciPy 1.0: Fundamental Algorithms for Scientific Computing in Python ». \textit{Nature Methods}, vol. 17, 2020.
\bibitem{ref31} Hunter, J. D. « Matplotlib: A 2D Graphics Environment ». \textit{Computing in Science \& Engineering}, vol. 9, n\textsuperscript{o} 3, 2007.
\bibitem{ref32} Streamlit. « Streamlit : bibliothèque Python open source pour applications de données ». Logiciel, version 1.57. \url{https://streamlit.io}
\bibitem{ref33} Chroma. « Chroma : base de données d'embeddings open source ». Logiciel, version 1.5. \url{https://www.trychroma.com}
\bibitem{ref34} Robertson, S. et Zaragoza, H. « The Probabilistic Relevance Framework: BM25 and Beyond ». \textit{Foundations and Trends in Information Retrieval}, vol. 3, n\textsuperscript{o} 4, 2009.
\bibitem{ref35} Brown, D. « Rank-BM25: A Collection of BM25 Algorithms in Python ». Dépôt GitHub, 2020. \url{https://github.com/dorianbrown/rank_bm25}
\bibitem{ref36} Ollama. « Ollama : exécution locale de modèles de langage ». Logiciel, version 0.34. \url{https://ollama.com}
\bibitem{ref37} W3C Web Incubator Community Group. « Web Speech API ». Spécification (Draft Community Group Report). \url{https://wicg.github.io/speech-api/}
\bibitem{ref38} Zaharia, M. et al. « Accelerating the Machine Learning Lifecycle with MLflow ». \textit{IEEE Data Engineering Bulletin}, vol. 41, n\textsuperscript{o} 4, 2018.
\bibitem{ref39} Krekel, H. et al. « pytest ». Logiciel, 2004. \url{https://pytest.org}
\bibitem{ref40} Jiang, A. Q. et al. « Mistral 7B ». \textit{arXiv:2310.06825}, 2023.
\bibitem{ref41} Laboratoire de Biologie, Faculté de Médecine et de Pharmacie de Sousse. « Images de microscopie de scratch assay annotées manuellement (97 images) ». Jeu de données interne, non publié, 2026.

\end{thebibliography}

\end{document}
